\documentclass[a4paper,12pt]{article}

\usepackage[utf8]{inputenc}
\usepackage[T1]{fontenc}
\usepackage[french]{babel}
\usepackage{graphicx}
\usepackage{geometry}
\usepackage[hidelinks]{hyperref}

% Illustrations et schémas
\usepackage{graphicx}
\usepackage{xcolor}
\usepackage{tikz}
\usepackage{pgfplots}
\usepackage{caption}
\usepackage{subcaption}
\usetikzlibrary{arrows.meta,positioning,calc,shapes,fit,backgrounds}
\pgfplotsset{compat=1.18}

% Tableaux (\toprule, \midrule, \bottomrule)
\usepackage{booktabs}

% Extraits de code
\usepackage{listings}
\definecolor{codebg}{RGB}{247,248,250}
\definecolor{codepurple}{RGB}{105,45,120}
\definecolor{codegreen}{RGB}{35,115,75}
\definecolor{codeblue}{RGB}{35,75,145}
\definecolor{lightblue}{RGB}{232,241,250}
\definecolor{lightgreen}{RGB}{233,245,236}
\definecolor{lightorange}{RGB}{255,243,224}

\lstdefinestyle{java}{
  language=Java,
  basicstyle=\ttfamily\small,
  keywordstyle=\color{codeblue}\bfseries,
  commentstyle=\color{codegreen}\itshape,
  stringstyle=\color{codepurple},
  backgroundcolor=\color{codebg},
  frame=single,
  rulecolor=\color{black!25},
  breaklines=true,
  columns=fullflexible,
  showstringspaces=false,
  tabsize=2,
  xleftmargin=0.5em,
  xrightmargin=0.5em,
  inputencoding=utf8,
  extendedchars=true,
  literate={é}{{\'e}}1 {è}{{\`e}}1 {ê}{{\^e}}1 {à}{{\`a}}1 {ç}{{\c{c}}}1 {É}{{\'E}}1 {û}{{\^u}}1 {î}{{\^i}}1
}
\lstset{style=java}

\geometry{margin=1.8cm}

\usepackage{setspace}
\onehalfspacing

% Repère visuel pour les passages qu'il me reste à écrire
\newcommand{\acompleter}[1]{\par\noindent\fbox{\parbox{0.96\linewidth}{\textcolor{red}{\textbf{REMARQUE :} #1}}}\par\medskip}

\begin{document}

\begin{titlepage}
    \centering

    % Logo
    \includegraphics[width=0.4\textwidth]{logoUM.png}\par\vspace{1.5cm}

    % Université
    {\scshape\LARGE Université de Montpellier \par}
    \vspace{0.5cm}
    {\scshape\Large Faculté des Sciences \par}
    \vspace{1cm}

    % Titre
    {\huge\bfseries HAI708I -- Entrepôts de données et Big-Data\par}
    \vspace{0.3cm}
    {\Large Projet EDBD \par}
    \vspace{2cm}

    % Auteurs
    {\Large
    Chanaud Léo
    \par}
    \vspace{1cm}

    % Num_étu
    {\Large
    Khalil Elhjaoui
    \par}
    \vspace{1cm}

    % Niveau
    {\large \textbf{MASTER IASD}\par}

    \vfill

    % Date
    {\large 2026 - 2027 \par}

\end{titlepage}

\small
\setstretch{1.1}
\sloppy

\tableofcontents

\newpage

% ============================================================
%  PARTIE 1 -- Analyse des besoins métier
% ============================================================

\section{Analyse des besoins métier}

% ------------------------------------------------------------
\subsection{Sources documentaires}
% ------------------------------------------------------------

Notre analyse s'appuie sur quatre sources, dont aucune n'est générée par une IA.

\begin{enumerate}
  \item \textbf{ACPM} (Alliance pour les chiffres de la presse et des médias) --
        déclarations de diffusion de la presse quotidienne nationale.
  \begin{itemize}
    \item Classement du second semestre 2025 :\\
          {\footnotesize\url{https://www.acpm.fr/classements/onenext-pqn}}
    \item Résultats 2021 (archive) :\\
          {\footnotesize\url{https://archive.wikiwix.com/cache/?url=https://www.acpm.fr/Les-chiffres/Audience-Presse/Resultats-par-etudes/OneNext2/Presse-Quotidienne-Nationale}}
  \end{itemize}

  \item \textbf{Reuters Institute} -- \emph{Digital News Report 2025} :\\
        {\footnotesize\url{https://reutersinstitute.politics.ox.ac.uk/sites/default/files/2025-06/Digital_News-Report_2025.pdf}}

  \item \textbf{Google Cloud} -- étude de cas sur le Groupe Le Monde :\\
        {\footnotesize\url{https://cloud.google.com/customers/le-monde-publicite}}

  \item \textbf{Le Monde} -- \og Un modèle économique basé sur les lecteurs
        et les abonnés \fg{} (2021) :\\
        {\footnotesize\url{https://www.lemonde.fr/le-monde-et-vous/article/2021/01/26/les-revenus-du-monde-des-sources-diversifiees_6067680_6065879.html}}
\end{enumerate}

% ------------------------------------------------------------
\subsection{Analyse du cas}
% ------------------------------------------------------------

\noindent\emph{Préambule : cette réponse, bien que rédigée par nos soins,
a été remise en page avec l'aide d'une IA afin d'être plus agréable à lire.}

\medskip

Ces sources font ressortir trois points essentiels pour la presse en ligne,
et en particulier pour un grand quotidien comme \emph{Le Monde}.

\paragraph{Répartition des revenus.}
Selon l'article du \emph{Monde} (2021), les revenus du groupe proviennent :

\begin{center}
\begin{tabular}{@{}lr@{}}
  \toprule
  \textbf{Source de revenu} & \textbf{Part} \\
  \midrule
  Abonnements numériques et papier & 54\,\% \\
  Recettes publicitaires            & 20\,\% \\
  Vente papier au numéro            & 14\,\% \\
  \bottomrule
\end{tabular}
\end{center}

Cette répartition montre une transition majeure : le modèle économique repose
désormais d'abord sur la \textbf{fidélisation des lecteurs} (les abonnés),
et non plus sur la vente ponctuelle ou sur la publicité seule.

\paragraph{Positionnement et modèle \og freemium \fg{}.}
Les données de l'ACPM confirment cette tendance : \emph{Le Monde} est un
leader de la presse quotidienne nationale, porté par une forte croissance de
sa diffusion numérique, au détriment du papier. Pour soutenir cette position,
le journal propose un modèle \og freemium \fg{} :

\begin{itemize}
  \item les contenus d'actualité sont en \textbf{accès libre}, pour générer du trafic ;
  \item les contenus à forte valeur ajoutée (enquêtes, dossiers, analyses)
        sont placés derrière un \textbf{mur payant} (\emph{paywall}),
        accessible uniquement sur abonnement.
\end{itemize}

\paragraph{Objectifs stratégiques et enjeux.}
Comme le souligne le rapport du Reuters Institute, convaincre les internautes
de payer pour de l'information en ligne reste un défi complexe. L'objectif
stratégique principal de l'entreprise est donc double :

\begin{enumerate}
  \item \textbf{convertir les lecteurs gratuits en abonnés numériques réguliers},
        afin de garantir son indépendance ;
  \item \textbf{optimiser le ciblage des espaces publicitaires} restants,
        comme l'illustre l'étude de cas Google Cloud.
\end{enumerate}

\paragraph{Besoins d'aide à la décision.}
Pour piloter cette double stratégie, la direction a besoin d'informations
précises. Elle doit pouvoir mesurer quels \textbf{articles ou rubriques} déclenchent le plus d'abonnements, identifier les \textbf{profils de lecteurs} qui risquent de résilier leur contrat et analyser la rentabilité des différents \textbf{canaux d'acquisition} (d'où provient le lecteur) et des \textbf{emplacements publicitaires}.

% ------------------------------------------------------------
\subsection{Actions et opérations à tracer}
% ------------------------------------------------------------

Pour un média comme \emph{Le Monde}, les événements qui créent de la valeur
se répartissent en trois processus principaux.

\begin{center}
\renewcommand{\arraystretch}{1.3}
\begin{tabular}{@{}p{0.04\linewidth}p{0.30\linewidth}p{0.56\linewidth}@{}}
  \toprule
  \textbf{\#} & \textbf{Action} & \textbf{Événement déclencheur} \\
  \midrule
  1 & \textbf{Souscription d'un abonnement} \newline
      {\footnotesize (Ventes)}
    & Un événement survient sur un contrat d'abonnement : souscription,
      renouvellement, changement d'offre ou résiliation. \\
  2 & \textbf{Consultation d'un article} \newline
      {\footnotesize (Web Analytics / Clickstream)}
    & Un lecteur charge la page d'un article. \\
  3 & \textbf{Affichage publicitaire} \newline
      {\footnotesize (Monétisation annexe)}
    & Une bannière publicitaire est chargée sur l'écran du lecteur. \\
  \bottomrule
\end{tabular}
\end{center}

% ------------------------------------------------------------
\subsection{Traitements analytiques et besoins métier}
% ------------------------------------------------------------

Pour chaque action, nous proposons des traitements (requêtes analytiques)
et indiquons le besoin métier auquel chacun répond. Les actions sont classés par ordre d'importance pour la question suivante.

% --- Action 1 ---
\subsubsection*{Action 1 -- Souscription d'un abonnement}

\begin{center}
\renewcommand{\arraystretch}{1.3}
\begin{tabular}{@{}p{0.45\linewidth}p{0.49\linewidth}@{}}
  \toprule
  \textbf{Traitement} & \textbf{Besoin métier} \\
  \midrule
  Nombre d’événements d’abonnement par type (souscription, renouvellement, changement d’offre, résiliation) et par mois.
    & Suivre le solde net d’abonnés (entrées moins sorties) et vérifier que la croissance n’est pas annulée par les départs \\
  \addlinespace
  Taux de désabonnement selon l'ancienneté du client et la région
  géographique.
    & Cibler les campagnes de fidélisation. \\
  \addlinespace
  Nombre de nouveaux abonnements et revenu moyen par abonné, selon le canal d’acquisition (paywall d’un article, newsletter, réseau social, moteur de recherche, partenaire) et la campagne marketing.
    & Répartir le budget marketing vers les canaux qui recrutent le plus d'abonnés (rentabilité) \\
  \bottomrule
\end{tabular}
\end{center}

% --- Action 2 ---
\subsubsection*{Action 2 -- Consultation d'un article}

\begin{center}
\renewcommand{\arraystretch}{1.3}
\begin{tabular}{@{}p{0.45\linewidth}p{0.49\linewidth}@{}}
  \toprule
  \textbf{Traitement} & \textbf{Besoin métier} \\
  \midrule
  Nombre de vues par rubrique et par support sur les 30 derniers jours.
    & Adapter l'interface utilisateur et orienter les choix éditoriaux
      selon les supports les plus utilisés. \\
  \addlinespace
  Taux de lecture complète des articles Premium \emph{vs} Gratuits.
    & Évaluer l'engagement des lecteurs face au contenu payant,
      pour justifier le prix de l'abonnement. \\
  \addlinespace
  Nombre d'articles consultés par un utilisateur anonyme avant qu'il
  ne quitte le site.
    & Ajuster l'algorithme du paywall (par exemple, bloquer la lecture
      au bout de $x$ articles pour inciter à l'abonnement). \\
  \bottomrule
\end{tabular}
\end{center}

\noindent\textit{Remarque :} le premier traitement peut être agrégé sur
différentes périodes (jour, semaine, mois, année) à l'aide des opérateurs
\texttt{CUBE} et \texttt{ROLLUP}.

% --- Action 3 ---
\subsubsection*{Action 3 -- Affichage publicitaire}

\begin{center}
\renewcommand{\arraystretch}{1.3}
\begin{tabular}{@{}p{0.45\linewidth}p{0.49\linewidth}@{}}
  \toprule
  \textbf{Traitement} & \textbf{Besoin métier} \\
  \midrule
  Revenu généré par emplacement publicitaire (Header, milieu d'article)
  et par jour.
    & Optimiser la disposition du site web pour maximiser les revenus
      sans nuire à la lecture. \\
  \addlinespace
  Taux de clic sur les publicités, croisé avec l'heure de la journée.
    & Ajuster le prix des espaces publicitaires vendus aux annonceurs
      selon les heures de pointe. \\
  \addlinespace
  Ratio entre le nombre de publicités affichées et le temps de lecture
  total d'une session.
    & Équilibrer monétisation et expérience utilisateur, pour éviter
      la fuite des lecteurs. \\
  \bottomrule
\end{tabular}
\end{center}

\bigskip

\textbf{Pour ces actions, on s'occupera de traiter dans le modèle détaillé l'action des abonnements, dans le modèle Snapshot on s'intéressera à l'action de consultation d'un article et dans le modèle updated-records on repartira sur les abonnements et on s'intéressera à l'historique (cycle de vie) qui en découle}.

\section{Conception du premier modèle détaillé}

D’après ce qui a été dit à la partie précédente, l’action que nous considérons la plus importante sont les Abonnements. Sa table de fait doit être au niveau le plus détaillé possible, on donne UNE LIGNE par évènement d'abonnement. On les modélise comme une souscription, chaque renouvellement facturé, chaque changement d’offre, résiliation ou réactivation sont des évènement. Ainsi on peut calculer par exemple le nombre actif d’abonnés, les groupes de personnes partageant la même caractéristique, le ARPU (Average Revenue Per User)...

% ============================================================
%  Schéma en étoile -- datamart détaillé des abonnements
%  Bibliothèques TikZ utilisées : positioning, calc, shapes
%  (déjà chargées dans le préambule du rapport)
%  Couleurs utilisées : lightblue, lightorange, lightgreen, codeblue
% ============================================================

\begin{figure}[htbp]
\centering
\begin{tikzpicture}[
  font=\footnotesize,
  dim/.style={
    rectangle split, rectangle split parts=2,
    rectangle split part fill={lightblue,white},
    draw=black!55, line width=0.6pt, rounded corners=2pt,
    text width=4.6cm, align=left, inner sep=4pt
  },
  fact/.style={
    rectangle split, rectangle split parts=4,
    rectangle split part fill={lightorange,white,white,lightgreen},
    draw=codeblue, line width=1.3pt, rounded corners=2pt,
    text width=5.4cm, align=left, inner sep=4pt
  },
  lien/.style={draw=black!55, line width=0.8pt},
  role/.style={font=\scriptsize\itshape, text=black!70}
]

% ---------------- Table des faits (au centre) ----------------
\node[fact] (F) at (0,0) {
  \textbf{FAIT\_EVENEMENT\_ABONNEMENT}\\
  {\scriptsize\itshape Grain : 1 événement d'abonnement}
  \nodepart{two}
  \textbf{Clés étrangères}\\
  date\_evenement\_sk (FK)\\
  date\_souscription\_sk (FK)\\
  heure\_sk (FK)\\
  lecteur\_sk (FK)\\
  offre\_sk (FK)\\
  promotion\_sk (FK)\\
  geo\_sk (FK)\\
  canal\_sk (FK)\\
  type\_evenement\_sk (FK)
  \nodepart{three}
  \textbf{Dimensions dégénérées}\\
  num\_contrat\\
  num\_transaction
  \nodepart{four}
  \textbf{Mesures}\\
  montant\_ttc \hfill [A]\\
  montant\_ht \hfill [A]\\
  montant\_remise \hfill [A]\\
  nb\_mois\_factures \hfill [A]\\
  delta\_abonnes ($+1$ / $0$ / $-1$) \hfill [A]\\
  revenu\_cumule\_abonne \hfill [SA]\\
  anciennete\_mois \hfill [NA]\\
  taux\_remise \hfill [NA]
};

% ---------------- Rangée du haut ----------------
\node[dim, anchor=south] (DATE) at ($(F.north)+(-5.9,1.4)$) {
  \textbf{DIM\_DATE}\\
  {\scriptsize\itshape Rôles : événement, souscription initiale}
  \nodepart{two}
  \textbf{date\_sk} (PK)\\
  date\_complete\\
  jour\_semaine\\
  est\_weekend\\
  semaine\_iso\\
  mois, libelle\_mois\\
  trimestre, semestre\\
  annee\\
  est\_jour\_ferie\\
  periode\_commerciale
};

\node[dim, anchor=south] (LEC) at ($(F.north)+(0,1.4)$) {
  \textbf{DIM\_LECTEUR}\\
  {\scriptsize\itshape Profil de l'abonné}
  \nodepart{two}
  \textbf{lecteur\_sk} (PK)\\
  id\_compte\_source\\
  tranche\_age\\
  categorie\_socio\_pro\\
  est\_etudiant\\
  temps\_visite\_jour (tranche)\\
  temps\_visite\_semaine (tranche)\\
  rubriques\_preferees\\
  niveau\_engagement\\
  date\_creation\_compte\\
  debut\_validite, fin\_validite
};

\node[dim, anchor=south] (OFF) at ($(F.north)+(5.9,1.4)$) {
  \textbf{DIM\_OFFRE}\\
  {\scriptsize\itshape Formule souscrite}
  \nodepart{two}
  \textbf{offre\_sk} (PK)\\
  code\_offre, libelle\_offre\\
  est\_premium\\
  support (numérique, papier)\\
  periodicite (mois, année)\\
  prix\_catalogue\_ttc\\
  duree\_engagement\_mois\\
  inclusions (Courrier int....)\\
  date\_lancement\\
  est\_active
};

% ---------------- Rangée du milieu ----------------
\node[dim] (HEU) at (-5.9,0) {
  \textbf{DIM\_HEURE}\\
  {\scriptsize\itshape Moment de la journée}
  \nodepart{two}
  \textbf{heure\_sk} (PK)\\
  heure (0 à 23)\\
  minute\\
  tranche\_horaire (matin, midi...)\\
  est\_heure\_de\_pointe\\
  apres\_envoi\_newsletter
};

\node[dim] (PRO) at (5.9,0) {
  \textbf{DIM\_PROMOTION}\\
  {\scriptsize\itshape Remise temporaire}
  \nodepart{two}
  \textbf{promotion\_sk} (PK)\\
  code\_promo\\
  type\_remise\\
  valeur\_remise\\
  duree\_promo\_mois\\
  public\_cible\\
  date\_debut, date\_fin\\
  evenement (Black Friday...)\\
  {\scriptsize\itshape 0 = aucune promotion}
};

% ---------------- Rangée du bas ----------------
\node[dim, anchor=north] (GEO) at ($(F.south)+(-5.9,-1.4)$) {
  \textbf{DIM\_GEOGRAPHIE}\\
  {\scriptsize\itshape Du plus large au plus fin}
  \nodepart{two}
  \textbf{geo\_sk} (PK)\\
  pays\\
  zone (Europe, Afrique...)\\
  region\\
  departement\\
  ville\\
  commune\\
  code\_postal\\
  code\_insee
};

\node[dim, anchor=north] (CAN) at ($(F.south)+(0,-1.4)$) {
  \textbf{DIM\_CANAL}\\
  {\scriptsize\itshape Par où l'abonné est arrivé}
  \nodepart{two}
  \textbf{canal\_sk} (PK)\\
  libelle\_canal\\
  famille (paywall, appli...)\\
  plateforme (web, iOS...)\\
  type\_appareil\\
  est\_payant\\
  source\_utm, medium\_utm\\
  nom\_campagne\\
  partenaire\\
  cout\_unitaire\_estime
};

\node[dim, anchor=north] (TYP) at ($(F.south)+(5.9,-1.4)$) {
  \textbf{DIM\_TYPE\_EVENEMENT}\\
  {\scriptsize\itshape Nature et cause de l'événement}
  \nodepart{two}
  \textbf{type\_evenement\_sk} (PK)\\
  type\_evenement\\
  sens (entrée, maintien, sortie)\\
  motif\_resiliation\\
  initiateur (abonné, système)\\
  moyen\_paiement\\
  statut\_paiement (refusé...)\\
  offre\_retention\_proposee\\
  offre\_retention\_acceptee
};

% ---------------- Liens ----------------
% DIM_DATE : deux liens, un par rôle
\draw[lien] ($(DATE.south)+(-0.3,0)$) -- ($(F.north)+(-2.2,0)$)
  node[role, pos=0.45, left=2pt] {2 rôles};
\draw[lien] ($(DATE.south)+(0.3,0)$)  -- ($(F.north)+(-1.6,0)$);

\draw[lien] (LEC.south) -- (F.north);
\draw[lien] (OFF.south) -- ($(F.north)+(1.9,0)$);

\draw[lien] (HEU.east) -- (HEU.east -| F.west);
\draw[lien] (PRO.west) -- (PRO.west -| F.east);

\draw[lien] (GEO.north) -- ($(F.south)+(-1.9,0)$);
\draw[lien] (CAN.north) -- (F.south);
\draw[lien] (TYP.north) -- ($(F.south)+(1.9,0)$);

\end{tikzpicture}

\caption{Schéma en étoile du datamart détaillé des abonnements.
[A] mesure additive, [SA] semi-additive (bloquée par DIM\_DATE),
[NA] non additive.}
\label{fig:etoile-abonnements}
\end{figure}

\newpage

\begin{description}

  \item[DIM\_LECTEUR] Un abonnement est souscrit par un lecteur qui fréquente
  le site et qui a décidé de s'abonner parce qu'il apprécie les articles.
  Plutôt que des attributs d'identité (nom, prénom), qui n'apportent rien à
  l'analyse, on y stocke son \textbf{profil de lecture} : temps de visite
  quotidien et hebdomadaire, rubriques préférées, niveau d'engagement,
  catégorie socio-professionnelle, tranche d'âge.

  \item[DIM\_OFFRE] Le journal propose plusieurs formules, premium ou non.
  Cette dimension permet de comparer les offres et de suivre les
  \textbf{montées ou descentes en gamme}. Elle décrit aussi les inclusions
  supplémentaires (\emph{Courrier international} ou un autre titre partenaire)
  et la date de lancement de l'offre.

  \item[DIM\_GEOGRAPHIE] Elle localise l'abonné à plusieurs niveaux, du plus
  large au plus fin : pays $>$ région $>$ département $>$ ville $>$ commune.
  Elle sert aux \textbf{analyses régionales} et permet de passer d'un niveau
  à l'autre (\emph{drill-down} et \emph{roll-up}).

  \item[DIM\_CANAL] Elle indique \textbf{par où l'abonné est arrivé} :
  paywall d'un article, application, newsletter, réseau social, partenaire
  ou campagne marketing. Elle mesure quels canaux recrutent le mieux,
  et à quel coût.

  \item[DIM\_PROMOTION] Elle décrit la remise appliquée (type, valeur, durée,
  public visé). Elle permet de vérifier si les abonnés venus grâce à une
  promotion \textbf{restent ou partent à la fin de l'offre}.
  On pourrait l'intégrer à \texttt{DIM\_OFFRE}, mais une offre dure plusieurs
  années alors qu'une promotion ne dure parfois que quelques semaines
  (Black Friday, soldes, rentrée). Les fusionner ferait exploser le nombre
  de lignes, avec une ligne par combinaison offre $\times$ promotion.
  De plus, la promotion concerne l'événement et non l'offre : au
  renouvellement suivant, l'abonné garde la même offre mais n'a plus de
  promotion.

  \item[DIM\_DATE] Elle situe chaque événement dans le temps (jour, mois,
  trimestre, jours fériés, périodes commerciales). Elle est utilisée dans
  \textbf{deux rôles} : la date de l'événement, et la date de souscription
  initiale qui définit la \emph{cohorte} de l'abonné.

  \item[DIM\_HEURE] Elle situe l'événement dans la journée (heure, tranche
  horaire, heure de pointe). Elle permet par exemple de voir si les
  abonnements arrivent juste après l'envoi de la newsletter du matin.
  Elle est séparée de \texttt{DIM\_DATE} pour garder les deux tables petites :
  365 jours + 1\,440 minutes, au lieu de 525\,600 lignes par an si l'heure
  était dans la date.

  \item[DIM\_TYPE\_EVENEMENT] Elle qualifie l'événement : son type, son sens
  (entrée, maintien, sortie), le motif de résiliation, l'initiateur et le
  statut du paiement. Elle sert à comprendre \textbf{pourquoi les abonnés
  entrent ou sortent}. Par exemple, elle distingue un abonné qui résilie
  volontairement d'un abonné dont la résiliation est déclenchée par le
  système après plusieurs paiements refusés.

\end{description}

% ------------------------------------------------------------
\subsection{Mesures}
% ------------------------------------------------------------

La table des faits porte huit mesures. Le type d'additivité indique sur
quelles dimensions on peut les sommer.

\medskip

\noindent\textit{Remarque :} nous n'avons pas retenu de mesures du type
\texttt{nb\_souscriptions} ou \texttt{nb\_resiliations}. Avec un grain d'un
événement par ligne, elles vaudraient toujours 0 ou 1 et répéteraient
l'information déjà portée par \texttt{DIM\_TYPE\_EVENEMENT}. Ces comptages
s'obtiennent par un \texttt{COUNT(*)} filtré sur le type d'événement, et le
solde net par \texttt{SUM(delta\_abonnes)}.

\begin{center}
\renewcommand{\arraystretch}{1.3}
\begin{tabular}{@{}p{0.24\linewidth}p{0.40\linewidth}p{0.30\linewidth}@{}}
  \toprule
  \textbf{Mesure} & \textbf{Définition} & \textbf{Type} \\
  \midrule
  montant\_ttc & Montant payé lors de l'événement (TTC). & Additive \\
  montant\_ht & Montant hors TVA (TVA presse de 2,1\,\%). & Additive \\
  montant\_remise & Prix catalogue moins montant payé. & Additive \\
  nb\_mois\_factures & Mois d'abonnement couverts par le paiement. & Additive \\
  delta\_abonnes & Variation du nombre d'abonnés ($+1$, $0$ ou $-1$). & Additive \\
  revenu\_cumule\_abonne & Total payé par l'abonné depuis sa première souscription. & Semi-additive, bloquée par \texttt{DIM\_DATE} \\
  anciennete\_mois & Mois écoulés depuis la souscription initiale. & Non additive \\
  taux\_remise & Part de la remise dans le prix catalogue. & Non additive \\
  \bottomrule
\end{tabular}
\end{center}

% ------------------------------------------------------------
\subsection{Précisions sur les attributs et mesures complexes}
% ------------------------------------------------------------

\begin{itemize}

  \item \textbf{delta\_abonnes} : vaut $+1$ quand un abonnement commence
  (souscription ou réactivation), $0$ quand l'abonné reste (renouvellement ou
  changement d'offre) et $-1$ lors d'une résiliation.
  \begin{itemize}
    \item Sa somme sur un mois donne le \emph{solde net} d'abonnés du mois.
    \item Sa somme sur tout l'historique donne le \emph{nombre d'abonnés actifs}.
  \end{itemize}

  \item \textbf{revenu\_cumule\_abonne} : total déjà payé par l'abonné,
  événement compris. On ne peut pas additionner deux valeurs du même abonné
  à deux dates différentes, car la seconde contient déjà la première.
  \begin{itemize}
    \item Exemple : un abonné paie 1\,€ en janvier puis 11,99\,€ en février.
          Les valeurs cumulées sont 1,00 puis 12,99\,€. Les additionner
          donnerait 13,99\,€, alors qu'il a payé 12,99\,€.
    \item Il faut prendre la dernière valeur de chaque abonné sur la période,
          puis sommer entre abonnés. C'est ce qui rend la mesure
          \emph{semi-additive} par rapport à la date.
  \end{itemize}

  \item \textbf{anciennete\_mois} et \textbf{taux\_remise} : additionner des
  anciennetés ou des pourcentages n'a pas de sens. On utilise la moyenne pour
  l'ancienneté. Le taux de remise se recalcule à partir des montants :
  \texttt{SUM(montant\_remise) / (SUM(montant\_ttc) + SUM(montant\_remise))}.

  \item \textbf{date\_evenement\_sk} et \textbf{date\_souscription\_sk} :
  deux clés vers la même table \texttt{DIM\_DATE}. La première situe
  l'événement, la seconde rattache l'abonné à sa \emph{cohorte}
  (par exemple, tous les abonnés de janvier 2026).

  \item \textbf{temps\_visite\_jour} et \textbf{temps\_visite\_semaine}
  (\texttt{DIM\_LECTEUR}) : stockés en \emph{tranches} (\og moins de 5 min \fg,
  \og 5 à 15 min \fg, etc.) recalculées chaque mois. Une valeur exacte
  changerait tous les jours et créerait une nouvelle version du lecteur à
  chaque variation, puisque la dimension est historisée.

  \item \textbf{niveau\_engagement} (\texttt{DIM\_LECTEUR}) : classe le lecteur
  selon sa fréquence de lecture (\og Quotidien \fg, \og Hebdomadaire \fg,
  \og Occasionnel \fg, \og Inactif \fg), recalculée chaque mois comme les
  temps de visite. Un passage de \og Quotidien \fg{} à \og Occasionnel \fg{}
  signale un décrochage, souvent annonciateur d'une résiliation.

  \item \textbf{debut\_validite} et \textbf{fin\_validite}
  (\texttt{DIM\_LECTEUR}) : historisation de type SCD~2. Quand le profil d'un
  lecteur change (par exemple, il n'est plus étudiant), on crée une nouvelle
  ligne au lieu d'écraser l'ancienne. Chaque événement reste ainsi rattaché
  au profil du lecteur \emph{au moment} de l'événement.

  \item \textbf{promotion\_sk = 0} (\texttt{DIM\_PROMOTION}) : ligne spéciale
  \og Aucune promotion \fg. Les événements sans remise y sont rattachés,
  ce qui évite des clés vides dans la table des faits.

  \item \textbf{sens} (\texttt{DIM\_TYPE\_EVENEMENT}) : regroupe les types
  d'événement en trois familles (entrée, maintien, sortie). Il sert à
  calculer \texttt{delta\_abonnes} lors du chargement.

\end{itemize}

\subsection{Répondre au besoin principal}

Oui, on peut répondre au besoin principal en s'appuyant sur notre Partie~1 : \textbf{convertir les lecteurs en abonnés et limiter les départs}. Comme chaque entrée et sortie d'abonné c'est une ligne de la table des faits, chaque traitement de l'action.

\begin{center}
\renewcommand{\arraystretch}{1.3}
\begin{tabular}{@{}p{0.38\linewidth}p{0.56\linewidth}@{}}
  \toprule
  \textbf{Traitement (Partie 1)} & \textbf{Comment le modèle y répond} \\
  \midrule
  Nombre d'événements d'abonnement par type et par mois.
    & \texttt{COUNT(*)} regroupé par \texttt{DIM\_TYPE\_EVENEMENT.type\_evenement}
      et \texttt{DIM\_DATE.mois}. Le solde net du mois s'obtient par
      \texttt{SUM(delta\_abonnes)}. \\
  \addlinespace
  Taux de désabonnement selon l'ancienneté et la région.
    & Les résiliations (\texttt{type\_evenement = 'Résiliation'}) sont
      regroupées par tranche de \texttt{anciennete\_mois} et par
      \texttt{DIM\_GEOGRAPHIE.region}. On les divise par le nombre d'abonnés
      actifs de la région, donné par \texttt{SUM(delta\_abonnes)} sur tout
      l'historique. \\
  \addlinespace
  Nouveaux abonnements et revenu moyen par canal et par campagne.
    & Les événements d'entrée (\texttt{sens = 'Entrée'}) sont comptés par
      \texttt{DIM\_CANAL.famille} et \texttt{nom\_campagne}. Le revenu moyen
      par abonné vaut \texttt{SUM(montant\_ttc) / COUNT(DISTINCT lecteur\_sk)}. \\
  \bottomrule
\end{tabular}
\end{center}

Les autres dimensions affinent ces analyses :
\begin{itemize}
  \item \texttt{DIM\_PROMOTION} indique si les abonnés recrutés grâce à une
        remise résilient à la fin de celle-ci ;
  \item \texttt{DIM\_HEURE} montre à quels moments de la journée les
        abonnements sont pris, par exemple après l'envoi de la newsletter ;
  \item le second rôle de \texttt{DIM\_DATE} (date de souscription)
        regroupe les abonnés par \emph{cohorte} pour suivre leur fidélité
        dans le temps.
\end{itemize}

\paragraph{Exemple de requête.}
La requête suivante répond au premier traitement. Grâce à \texttt{ROLLUP},
elle donne à la fois le détail par type d'événement et le total de chaque
mois. Sur les lignes de total (où \texttt{type\_evenement} est vide), la
colonne \texttt{solde\_net} indique directement si le journal a gagné
ou perdu des abonnés ce mois-là.

\begin{lstlisting}[language=SQL]
SELECT d.mois,
       t.type_evenement,
       COUNT(*)             AS nb_evenements,
       SUM(f.delta_abonnes) AS solde_net
FROM fait_evenement_abonnement f
JOIN dim_date d           ON f.date_evenement_sk = d.date_sk
JOIN dim_type_evenement t ON f.type_evenement_sk = t.type_evenement_sk
WHERE d.annee = 2026
GROUP BY ROLLUP (d.mois, t.type_evenement)
ORDER BY d.mois, t.type_evenement;
\end{lstlisting}

Si le solde net d'un mois est négatif alors que les souscriptions sont
nombreuses, la direction sait que le problème vient des départs et non du
recrutement. Elle peut alors déplacer son effort vers la fidélisation. Bien évidemment, ce modèle doit avoir des limites sur pleins d'aspects, mais c'est celui qui nous parait le plus cohérent. Petite justification aussi pour la requête ci-dessus : prendre d.annee = 2026 peut paraître non-analytique, mais dans Le Monde, les années peuvent vraiment évoluer, surtout dans notre société actuelle, donc il est intéressant de voir que la façon de consommer l'information nous pousse à analyser sur qu'une seulle année.

\subsection{Exemple d'instance}

On prend une instance, donc pour moi ça signifie par exemple deux abonnés fictifs. Une étudiante lyonnaise visite le site, et au bout de quelques jours, décide de prendre l’abonnement avec une promotion de 1€ puis renouvelle. Un cadre rennais résilie à la fin de sa deuxième année. Les prix des offres sont fictifs et seules les colonnes principales de chaque dim seront montrées.

Les tables de cette instance (dimensions et table des faits) sont données
en annexe~\ref{annexe:instance-detaille}.

\medskip

\noindent Lecture de l'instance :
\begin{itemize}
  \item \textbf{F1} : souscription avec promotion, l'étudiante paie 1\,€
        au lieu de 11,99\,€ (remise de 10,99\,€). La promotion est
        rattachée à l'événement : en \textbf{F2}, la même offre est
        renouvelée sans promotion (\texttt{promotion\_sk = 0}),
        automatiquement à 00:05 par le système.
  \item \textbf{F3} : la résiliation ne génère aucun paiement. Le revenu
        cumulé (398\,€) correspond aux deux années payées à 199\,€.
        Le canal reste celui d'acquisition, ce qui permet de mesurer la
        fidélité des abonnés selon leur canal d'origine.
  \item \textbf{Contrôles} : \texttt{SUM(montant\_ttc)} pour l'abonnée 1001
        donne bien 12,99\,€ payés. \texttt{SUM(delta\_abonnes)} sur janvier
        2026 vaut $+1 - 1 = 0$ (une entrée, une sortie). En revanche,
        additionner \texttt{revenu\_cumule\_abonne} de F1 et F2 donnerait
        13,99\,€ au lieu de 12,99\,€ : c'est la semi-additivité sur la date.
\end{itemize}

D’après les sources, le nombre d’abonnés fin 2025 sont d’environ 700k abonnés. On prends tout ces abonnés et on renouvelle tout les mois pendant 12, donc 12 x 700k = 8.4M de lignes au minimum, après peut être certains sont prélevés tout les ans plutôt donc ça fait descendre mais au pire des cas c’est environ 8 millions, sachant qu’on a aussi les souscriptions, les réactivations, les résiliations, les paiements refusés et les changements d’offres. Même si c’est des minorités pour un gros nombre d’abonnés c’est significatif. D’après Internet, un fichier sheets peut contenir au max 1M de lignes, donc pas gérable, surtout si on analyse sur plusieurs années pour les grp de personnes (cohortes). Imaginons sur 3 ans l’analyse devient 20M. Même niveau sécurite / RGPD, alimentation ou historique, sheets est un peu limité.

% ============================================================
%  PARTIE 3 -- Conception de deux modèles moins détaillés
% ============================================================
\section{Conception de deux modèles moins détaillés}

L'objectif de cette partie est de concevoir deux datamarts complémentaires reposant sur des grains moins fins que le modèle transactionnel de la Partie 2 : un modèle \emph{snapshot} (instantané périodique) portant sur la consultation des articles, et un modèle \emph{updated-records} (instantané cumulatif) dédié au suivi du cycle de vie des contrats d'abonnement.

% ============================================================
%  3.1 Datamart Snapshot -- Consultation quotidienne d'articles
% ============================================================

\subsection{Modèle Snapshot : consultation quotidienne d'articles}

Parmi les actions prioritaires identifiées en Partie 1, la consultation d'articles se prête naturellement à une modélisation en \emph{snapshot} périodique. Un grand quotidien en ligne générant plusieurs millions de clics par jour, stocker chaque événement de lecture unitairement entraînerait une volumétrie disproportionnée pour les requêtes courantes de pilotage éditorial.

Nous définissons donc la table de faits \texttt{FAITS\_CONSULTATION\_QUOTIDIENNE}, dont le grain correspond au \textbf{bilan quotidien des consultations d'un article, pour une source de trafic, un support technique et un profil d'audience donnés}. Seules les combinaisons ayant enregistré au moins une vue dans la journée génèrent une ligne à minuit.

% ---------------- Schéma en étoile TikZ : Snapshot ----------------
\begin{figure}[htbp]
\centering
\begin{tikzpicture}[
  font=\footnotesize,
  dim/.style={
    rectangle split, rectangle split parts=2,
    rectangle split part fill={lightblue,white},
    draw=black!55, line width=0.6pt, rounded corners=2pt,
    text width=4.6cm, align=left, inner sep=4pt
  },
  fact/.style={
    rectangle split, rectangle split parts=3,
    rectangle split part fill={lightorange,white,lightgreen},
    draw=codeblue, line width=1.3pt, rounded corners=2pt,
    text width=5.6cm, align=left, inner sep=4pt
  },
  lien/.style={draw=black!55, line width=0.8pt}
]

% Table des faits (au centre)
\node[fact] (FS) at (0,0) {
  \textbf{FAITS\_CONSULTATION\_QUOTIDIENNE}\\
  {\scriptsize\itshape Grain : jour $\times$ article $\times$ source $\times$ support $\times$ profil}
  \nodepart{two}
  \textbf{Clés étrangères}\\
  date\_sk (FK)\\
  article\_sk (FK)\\
  source\_sk (FK)\\
  support\_sk (FK)\\
  profil\_sk (FK)
  \nodepart{three}
  \textbf{Mesures}\\
  nb\_vues\_totales \hfill [A]\\
  nb\_lectures\_completes \hfill [A]\\
  temps\_lecture\_total\_sec \hfill [A]\\
  nb\_commentaires\_cumules \hfill [SA]\\
  temps\_lecture\_moyen \hfill [NA]\\
  taux\_completion \hfill [NA]
};

% Rangée du haut
\node[dim, anchor=south] (SDATE) at ($(FS.north)+(-3.4,1.2)$) {
  \textbf{DIM\_DATE}\\
  {\scriptsize\itshape Date de l'instantané quotidien}
  \nodepart{two}
  \textbf{date\_sk} (PK)\\
  date\_complete\\
  jour\_mois, jour\_semaine\\
  mois, libelle\_mois\\
  trimestre, annee\\
  est\_weekend\\
  est\_jour\_ferie\\
  periode\_scolaire
};

\node[dim, anchor=south] (SART) at ($(FS.north)+(3.4,1.2)$) {
  \textbf{DIM\_ARTICLE}\\
  {\scriptsize\itshape Contenu éditorial consulté}
  \nodepart{two}
  \textbf{article\_sk} (PK)\\
  id\_article\_source\\
  titre\_article\\
  rubrique\_principale\\
  sous\_rubrique\\
  format\_contenu (texte, vidéo...)\\
  type\_acces (Premium, Gratuit)\\
  auteur\_principal\\
  date\_publication\\
  longueur\_mots
};

% Rangée du milieu / bas
\node[dim] (SSRC) at (-6.0,-0.3) {
  \textbf{DIM\_SOURCE\_TRAFIC}\\
  {\scriptsize\itshape Provenance du clic}
  \nodepart{two}
  \textbf{source\_sk} (PK)\\
  nom\_source\\
  categorie\_source\\
  est\_trafic\_interne\\
  medium\_acquisition\\
  type\_emplacement\\
  est\_recommandation\_algo\\
  referent\_externe\_domaine\\
  nom\_campagne\_marketing\\
  position\_lien\_clique
};

\node[dim] (SSUP) at (6.0,-0.3) {
  \textbf{DIM\_SUPPORT}\\
  {\scriptsize\itshape Outil technique de lecture}
  \nodepart{two}
  \textbf{support\_sk} (PK)\\
  famille\_appareil\\
  systeme\_exploitation\\
  version\_os\\
  navigateur\\
  type\_application\\
  version\_application\\
  resolution\_ecran\\
  est\_mode\_sombre\_actif\\
  categorie\_reseau
};

\node[dim, anchor=north] (SSEG) at ($(FS.south)+(0,-1.2)$) {
  \textbf{DIM\_PROFIL\_LECTEUR}\\
  {\scriptsize\itshape Profil d'audience agrégé}
  \nodepart{two}
  \textbf{profil\_sk} (PK)\\
  statut\_abonnement\\
  niveau\_engagement\\
  tranche\_age\_declaree\\
  categorie\_socio\_pro\\
  preference\_rubrique\\
  a\_accepte\_cookies\\
  utilise\_adblock\\
  anciennete\_profil\\
  type\_authentification
};

% Liens
\draw[lien] (SDATE.south) -- ($(FS.north)+(-1.6,0)$);
\draw[lien] (SART.south)  -- ($(FS.north)+(1.6,0)$);
\draw[lien] (SSRC.east)   -- (SSRC.east -| FS.west);
\draw[lien] (SSUP.west)   -- (SSUP.west -| FS.east);
\draw[lien] (SSEG.north)  -- (FS.south);

\end{tikzpicture}
\caption{Schéma en étoile du datamart \emph{snapshot} des consultations d'articles.
[A] mesure additive, [SA] semi-additive (bloquée par DIM\_DATE), [NA] non additive.}
\label{fig:etoile-snapshot}
\end{figure}

\subsubsection*{Justification des dimensions du modèle Snapshot}

Ces cinq dimensions s'imposent logiquement au regard de la nature d'un modèle \emph{Snapshot} et des besoins analytiques identifiés en Partie 1 :

\begin{description}
  \item[DIM\_PROFIL\_LECTEUR] Pour un instantané quotidien qui pré-agrège des millions de clics, il est incohérent de conserver le lecteur individuel. Cette dimension découle naturellement du grain agrégé en regroupant l'audience par \textbf{profils types} (statut d'abonnement, tranche d'âge déclarée, niveau d'engagement, présence d'un bloqueur de publicité).
  \item[DIM\_ARTICLE] Elle décrit le contenu éditorial (rubrique, sous-rubrique, auteur, format texte ou vidéo, longueur en mots) et surtout son \textbf{modèle d'accès} (\texttt{Premium} ou \texttt{Gratuit}), indispensable pour évaluer l'attractivité des articles payants.
  \item[DIM\_SOURCE\_TRAFIC] Plutôt que de multiplier des colonnes de compteurs dans la table de faits (vues depuis la Une, vues par recommandation), cette dimension isole la \textbf{provenance de la visite} (accès interne, moteur de recherche, réseau social, algorithme de recommandation, campagne marketing).
  \item[DIM\_SUPPORT] Elle caractérise l'environnement technique de lecture (mobile, ordinateur, système d'exploitation, application native ou navigateur web) afin d'adapter l'ergonomie aux appareils réellement utilisés.
  \item[DIM\_DATE] Dimension conforme partagée avec la Partie 2, elle date la photographie quotidienne prise à minuit et permet d'analyser la saisonnalité de l'audience (week-ends, vacances scolaires).
\end{description}

\subsubsection*{Mesures du modèle Snapshot}

La table \texttt{FAITS\_CONSULTATION\_QUOTIDIENNE} comporte six mesures résumant l'activité de lecture de la journée :

\begin{center}
\renewcommand{\arraystretch}{1.3}
\begin{tabular}{@{}p{0.28\linewidth}p{0.38\linewidth}p{0.28\linewidth}@{}}
  \toprule
  \textbf{Mesure} & \textbf{Définition métier} & \textbf{Type d'additivité} \\
  \midrule
  \texttt{nb\_vues\_totales} & Nombre total de chargements de la page de l'article pour cette combinaison précise. & Additive \\
  \texttt{nb\_lectures\_completes} & Nombre de fois où les lecteurs ont fait défiler (\emph{scroll}) l'article jusqu'à la fin. & Additive \\
  \texttt{temps\_lecture\_total\_sec} & Somme du temps passé (en secondes) par tous les lecteurs du profil sur cet article. & Additive \\
  \texttt{nb\_commentaires\_cumules} & Stock total de commentaires validés sous l'article à la fin de la journée (minuit). & Semi-additive, bloquée par \texttt{DIM\_DATE} \\
  \texttt{temps\_lecture\_moyen} & Temps moyen passé par vue (\texttt{temps\_lecture\_total\_sec / nb\_vues\_totales}). & Non additive \\
  \texttt{taux\_completion} & Part des vues ayant abouti à une lecture complète (\texttt{nb\_lectures\_completes / nb\_vues\_totales}). & Non additive \\
  \bottomrule
\end{tabular}
\end{center}

\noindent\textit{Précision sur la semi-additivité :} \texttt{nb\_commentaires\_cumules} représente un état (stock) en fin de journée. On peut additionner les commentaires de deux articles différents le même jour, mais additionner le cumul d'un même article du lundi (ex. 50) avec celui du mardi (ex. 55) donnerait un résultat erroné (105 au lieu de 55). Les mesures \texttt{temps\_lecture\_moyen} et \texttt{taux\_completion} étant des ratios, elles ne s'additionnent sur aucune dimension et se recalculent dynamiquement via les sommes de leurs numérateurs et dénominateurs.

\subsubsection*{Réponse aux besoins métier (Snapshot)}

Ce modèle \emph{Snapshot} répond directement aux enjeux stratégiques du journal tout en garantissant des performances d'affichage rapides pour l'aide à la décision. Sur le plan éditorial, le croisement de \texttt{DIM\_ARTICLE} avec \texttt{temps\_lecture\_total\_sec} et \texttt{nb\_lectures\_completes} indique immédiatement à la rédaction quels sujets, formats ou articles \emph{Premium} retiennent réellement l'attention des lecteurs (Traitement 2 de l'Action 2). Sur le plan marketing et technique, l'analyse combinée de \texttt{DIM\_SOURCE\_TRAFIC} et de \texttt{DIM\_SUPPORT} permet à la fois de savoir où investir pour attirer du trafic qualifié et de repérer rapidement un problème d'ergonomie si le taux de complétion chute sur un appareil spécifique (Traitement 1 de l'Action 2). Enfin, en pré-agrégeant les clics chaque soir à minuit par profil de lecteur, le modèle évite de parcourir des milliards d'événements unitaires lors de l'interrogation sur plusieurs mois.

\subsubsection*{Exemple d'instance (Snapshot)}

Pour tester ce modèle, l'instance ci-dessous illustre deux situations courantes. La première ligne représente la lecture d'un long article politique payant (\og Les enjeux des élections... \fg{}), lu attentivement par des abonnés fidèles sur l'application mobile iOS depuis la page d'accueil. La seconde ligne montre le visionnage d'une courte vidéo de sport gratuite (\og Résultats du marathon \fg{}), qui attire massivement des visiteurs occasionnels arrivant depuis Google Discover sur ordinateur.

Les tables de cette instance sont données en
annexe~\ref{annexe:instance-snapshot}.

% ============================================================
%  3.2 Datamart Updated-Records -- Cycle de vie des abonnements
% ============================================================

\newpage
\subsection{Modèle Updated-Records : cycle de vie des contrats d'abonnement}

Le modèle \emph{updated-records} (\emph{accumulating snapshot} chez Kimball) suit un processus métier structuré par un début, des jalons connus et une fin. Nous l'appliquons à l'Action 1 (\textbf{les abonnements}) afin d'analyser le parcours complet d'un contrat de sa souscription jusqu'à sa résiliation.

Dans la Partie 2, un contrat produisait une nouvelle ligne à chaque événement (soit au moins 12 lignes par an pour un abonné mensuel). Ici, la table \texttt{FAIT\_CYCLE\_ABONNEMENT} a pour grain \textbf{une seule ligne par contrat d'abonnement} (identifiée par la dimension dégénérée \texttt{num\_contrat}), créée à la souscription et \textbf{mise à jour à chaque étape franchie} ou paiement enregistré.

% ---------------- Schéma en étoile TikZ : Updated-Records ----------------
\begin{figure}[htbp]
\centering
\begin{tikzpicture}[
  font=\footnotesize,
  dim/.style={
    rectangle split, rectangle split parts=2,
    rectangle split part fill={lightblue,white},
    draw=black!55, line width=0.6pt, rounded corners=2pt,
    text width=4.5cm, align=left, inner sep=4pt
  },
  fact/.style={
    rectangle split, rectangle split parts=4,
    rectangle split part fill={lightorange,white,white,lightgreen},
    draw=codeblue, line width=1.3pt, rounded corners=2pt,
    text width=5.6cm, align=left, inner sep=4pt
  },
  lien/.style={draw=black!55, line width=0.8pt},
  role/.style={font=\scriptsize\itshape, text=black!70}
]

% Table des faits (au centre)
\node[fact] (FU) at (0,0) {
  \textbf{FAIT\_CYCLE\_ABONNEMENT}\\
  {\scriptsize\itshape Grain : 1 ligne par contrat (mise à jour)}
  \nodepart{two}
  \textbf{Clés étrangères}\\
  date\_souscription\_sk (FK)\\
  date\_fin\_promo\_sk (FK)\\
  date\_plein\_tarif\_sk (FK)\\
  date\_demande\_resil\_sk (FK)\\
  date\_fin\_contrat\_sk (FK)\\
  date\_derniere\_maj\_sk (FK)\\
  lecteur\_sk (FK)\\
  offre\_initiale\_sk, offre\_actuelle\_sk (FK)\\
  canal\_sk, promotion\_sk, geo\_sk (FK)\\
  statut\_sk (FK)
  \nodepart{three}
  \textbf{Dimension dégénérée}\\
  num\_contrat
  \nodepart{four}
  \textbf{Mesures (principales)}\\
  nb\_contrats ($=1$) \hfill [A]\\
  montant\_total\_paye\_ttc \hfill [A]\\
  montant\_total\_remise \hfill [A]\\
  nb\_mois\_payes, nb\_chgt\_offre \hfill [A]\\
  est\_converti\_plein\_tarif ($0/1$) \hfill [A]\\
  est\_resilie ($0/1$) \hfill [A]\\
  revenu\_mensuel\_actuel (MRR) \hfill [SA]\\
  jours\_engagement\_restants \hfill [SA]\\
  delai\_conversion\_jours \hfill [NA]\\
  duree\_vie\_jours, delai\_preavis \hfill [NA]
};

% Rangée du haut
\node[dim, anchor=south] (UDATE) at ($(FU.north)+(-5.6,1.2)$) {
  \textbf{DIM\_DATE}\\
  {\scriptsize\itshape 6 rôles de jalons temporels}
  \nodepart{two}
  \textbf{date\_sk} (PK)\\
  date\_complete\\
  mois, trimestre, annee\\
  {\scriptsize\itshape + ligne spéciale \texttt{99991231}}\\
  {\scriptsize\itshape \og Étape non atteinte \fg{}}
};

\node[dim, anchor=south] (ULEC) at ($(FU.north)+(0,1.2)$) {
  \textbf{DIM\_LECTEUR}\\
  {\scriptsize\itshape Conforme Partie 2 (SCD 2)}
  \nodepart{two}
  \textbf{lecteur\_sk} (PK)\\
  id\_compte\_source\\
  tranche\_age, CSP\\
  est\_etudiant\\
  profil\_engagement
};

\node[dim, anchor=south] (UOFF) at ($(FU.north)+(5.6,1.2)$) {
  \textbf{DIM\_OFFRE}\\
  {\scriptsize\itshape 2 rôles : initiale et actuelle}
  \nodepart{two}
  \textbf{offre\_sk} (PK)\\
  code\_offre, libelle\_offre\\
  gamme, support\\
  periodicite\_facturation\\
  prix\_catalogue\_ttc
};

% Rangée du milieu
\node[dim] (UCAN) at (-5.9,0) {
  \textbf{DIM\_CANAL}\\
  {\scriptsize\itshape Canal d'acquisition initial}
  \nodepart{two}
  \textbf{canal\_sk} (PK)\\
  libelle\_canal, famille\_canal\\
  plateforme, type\_appareil\\
  source\_utm, medium\_utm\\
  cout\_unitaire\_estime
};

\node[dim] (UPRO) at (5.9,0) {
  \textbf{DIM\_PROMOTION}\\
  {\scriptsize\itshape Promotion de souscription}
  \nodepart{two}
  \textbf{promotion\_sk} (PK)\\
  code\_promo, libelle\_promo\\
  type\_remise, valeur\_remise\\
  duree\_promo\_mois\\
  prix\_promo\_ttc
};

% Rangée du bas
\node[dim, anchor=north] (UGEO) at ($(FU.south)+(-3.5,-1.2)$) {
  \textbf{DIM\_GEOGRAPHIE}\\
  {\scriptsize\itshape Adresse de facturation}
  \nodepart{two}
  \textbf{geo\_sk} (PK)\\
  code\_postal, commune\\
  departement, region\\
  pays, zone\\
  est\_zone\_portage\_papier
};

\node[dim, anchor=north] (USTAT) at ($(FU.south)+(3.5,-1.2)$) {
  \textbf{DIM\_STATUT\_CONTRAT}\\
  {\scriptsize\itshape État courant du cycle de vie}
  \nodepart{two}
  \textbf{statut\_sk} (PK)\\
  statut (En promo, Plein tarif...)\\
  numero\_etape (1 à 5)\\
  est\_actif, est\_payant\_plein\_tarif\\
  type\_fin, motif\_fin\\
  initiateur\_fin, canal\_resiliation\\
  offre\_retention\_proposee / acceptee\\
  est\_reconquerable
};

% Liens
\draw[lien] ($(UDATE.south)+(-0.2,0)$) -- ($(FU.north)+(-2.1,0)$)
  node[role, pos=0.45, left=2pt] {6 rôles};
\draw[lien] ($(UDATE.south)+(0.2,0)$)  -- ($(FU.north)+(-1.6,0)$);

\draw[lien] (ULEC.south) -- (FU.north);

\draw[lien] ($(UOFF.south)+(-0.2,0)$) -- ($(FU.north)+(1.6,0)$);
\draw[lien] ($(UOFF.south)+(0.2,0)$)  -- ($(FU.north)+(2.1,0)$)
  node[role, pos=0.45, right=2pt] {2 rôles};

\draw[lien] (UCAN.east) -- (UCAN.east -| FU.west);
\draw[lien] (UPRO.west) -- (UPRO.west -| FU.east);

\draw[lien] (UGEO.north)  -- ($(FU.south)+(-1.6,0)$);
\draw[lien] (USTAT.north) -- ($(FU.south)+(1.6,0)$);

\end{tikzpicture}
\caption{Schéma en étoile du datamart \emph{updated-records} du cycle d'abonnement.}
\label{fig:etoile-updated}
\end{figure}

\subsubsection*{Dimensions et rôles multiples}

Ce modèle mobilise sept dimensions. Six d'entre elles (\texttt{DIM\_DATE}, \texttt{DIM\_LECTEUR}, \texttt{DIM\_OFFRE}, \texttt{DIM\_CANAL}, \texttt{DIM\_PROMOTION}, \texttt{DIM\_GEOGRAPHIE}) sont des \textbf{dimensions conformes} strictement partagées avec la Partie 2, garantissant la cohérence des analyses croisées. Deux adaptations majeures caractérisent ce schéma :
\begin{itemize}
  \item \textbf{Rôles multiples sur \texttt{DIM\_DATE} et \texttt{DIM\_OFFRE} :} La ligne de faits embarque six clés de dates correspondant aux étapes clés du contrat (\emph{souscription}, \emph{fin de promotion}, \emph{passage au plein tarif}, \emph{demande de résiliation}, \emph{fin effective du contrat} et \emph{dernière mise à jour}). Tant qu'une étape n'est pas franchie, sa clé pointe vers l'enregistrement spécial \texttt{99991231} (\og Étape non atteinte \fg{}). De même, \texttt{DIM\_OFFRE} joue deux rôles (\texttt{offre\_initiale\_sk} et \texttt{offre\_actuelle\_sk}) pour suivre immédiatement les montées ou descentes en gamme.
  \item \textbf{Nouvelle dimension \texttt{DIM\_STATUT\_CONTRAT} :} Elle remplace \texttt{DIM\_TYPE\_EVENEMENT} pour décrire l'état courant du contrat à l'aide de 11 attributs : \texttt{statut\_sk}, \texttt{statut} (\emph{En promotion, Plein tarif, Résiliation demandée, Suspendu, Résilié}), \texttt{numero\_etape} (1 à 5), \texttt{est\_actif}, \texttt{est\_payant\_plein\_tarif}, \texttt{type\_fin} (\emph{Volontaire, Involontaire}), \texttt{motif\_fin}, \texttt{initiateur\_fin}, \texttt{canal\_resiliation}, \texttt{offre\_retention\_proposee}, \texttt{offre\_retention\_acceptee} et \texttt{est\_reconquerable}.
\end{itemize}

\subsubsection*{Mesures du modèle Updated-Records}

La table \texttt{FAIT\_CYCLE\_ABONNEMENT} regroupe 13 mesures recalculées à chaque mise à jour de la ligne :

\begin{center}
\renewcommand{\arraystretch}{1.25}
\begin{tabular}{@{}p{0.28\linewidth}p{0.40\linewidth}p{0.26\linewidth}@{}}
  \toprule
  \textbf{Mesure} & \textbf{Définition} & \textbf{Type d'additivité} \\
  \midrule
  \texttt{nb\_contrats} (M1) & Vaut toujours 1 ; sert au comptage des contrats. & Additive \\
  \texttt{montant\_total\_paye\_ttc} (M2) & Total encaissé sur le contrat depuis la souscription. & Additive \\
  \texttt{montant\_total\_remise} (M3) & Cumul des remises accordées sur le contrat. & Additive \\
  \texttt{nb\_mois\_payes} (M4) & Nombre total de mois d'abonnement facturés. & Additive \\
  \texttt{nb\_paiements\_refuses} (M5) & Nombre d'échecs de prélèvement sur le contrat. & Additive \\
  \texttt{nb\_changements\_offre} (M6) & Nombre de changements de formule effectués. & Additive \\
  \texttt{est\_converti\_plein\_tarif} (M7) & Vaut 1 si le contrat a payé au moins une fois au plein tarif, 0 sinon. & Additive \\
  \texttt{est\_resilie} (M8) & Vaut 1 si le contrat est terminé, 0 sinon. & Additive \\
  \texttt{delai\_conversion\_jours} (M9) & Jours écoulés entre la souscription et le 1\textsuperscript{er} paiement plein tarif. & Non additive \\
  \texttt{duree\_vie\_jours} (M10) & Jours entre la souscription et la fin (ou la dernière mise à jour). & Non additive \\
  \texttt{delai\_preavis\_jours} (M11) & Jours entre la demande de résiliation et la fin effective. & Non additive \\
  \texttt{revenu\_mensuel\_actuel} (M12) & Montant mensuel payé aujourd'hui (MRR), 0 si résilié. & Semi-additive (\texttt{DIM\_DATE}) \\
  \texttt{jours\_engagement\_restants} (M13) & Jours avant la fin d'engagement à la date de mise à jour. & Semi-additive (\texttt{DIM\_DATE}) \\
  \bottomrule
\end{tabular}
\end{center}

\noindent\textit{Changement de nature d'additivité par rapport à la Partie 2 :} Dans le modèle transactionnel (Partie 2), le revenu cumulé était semi-additif car un même abonné apparaissait sur plusieurs lignes datées. Dans ce modèle \emph{updated-records}, chaque contrat n'occupe qu'une seule ligne écrasée à chaque mise à jour : \texttt{montant\_total\_paye\_ttc} (M2) n'est jamais dupliqué et devient pleinement \textbf{additif}. En revanche, \texttt{revenu\_mensuel\_actuel} (M12) et \texttt{jours\_engagement\_restants} (M13) sont des soldes à l'instant $t$ (semi-additifs sur la date) ; l'ETL rafraîchit chaque nuit toutes les lignes actives pour aligner leur \texttt{date\_derniere\_maj\_sk}.

\subsubsection*{Réponse aux besoins métier (Updated-Records)}

Ce modèle est le plus direct pour analyser la fidélisation, la valeur vie client et la vitesse de parcours d'un abonné. Comme toutes les dates d'étape figurent sur la même ligne, calculer le délai de conversion après une promotion ou le taux de départ avant le plein tarif ne nécessite aucune auto-jointure complexe entre événements.

\bigskip

Si l'offre \og 1\,€ le premier mois \fg{} recrute fortement mais affiche un fort taux de départ avant le premier paiement plein tarif, la direction peut immédiatement la remplacer par une promotion étalée sur trois mois. La limite de ce modèle est l'absence d'états intermédiaires (un abonné ayant changé trois fois d'offre ne conserve que l'offre initiale et l'offre actuelle), ce qui justifie sa complémentarité avec le modèle détaillé de la Partie 2.

\subsubsection*{Exemple d'instance (Updated-Records)}

Nous reprenons les deux abonnés fictifs de la Partie 2 (les dimensions \texttt{DIM\_LECTEUR}, \texttt{DIM\_OFFRE}, \texttt{DIM\_CANAL}, \texttt{DIM\_PROMOTION} et \texttt{DIM\_GEOGRAPHIE} restent identiques). Le contrat \texttt{CT-5001} (étudiante lyonnaise) est représenté à travers \textbf{trois états successifs} ($E_1, E_2, E_3$) de la même ligne physique réécrite au fil des mois, tandis que \texttt{CT-1187} (cadre rennais) est à l'état résilié.

Les tables de cette instance sont données en
annexe~\ref{annexe:instance-updated}.

Après le 15 mars 2026, la table physique ne contient que deux lignes ($E_3$ et \texttt{CT-1187}). La somme $\texttt{SUM(M2)} = 24{,}98 + 398{,}00 = 422{,}98\,\text{€}$ donne immédiatement le chiffre d'affaires encaissé sans double comptage.


% ============================================================
%  3.3 Volumétrie sur 12 mois et comparaison avec Excel (Q.20)
% ============================================================

\subsection{Estimation de la volumétrie sur 12 mois et pertinence face à un tableur}

Les deux datamarts présentent des profils de volumétrie opposés, basés sur les chiffres publics du \emph{Monde} (environ $670\,000$ abonnés numériques et $181{,}1$ millions de visites mensuelles, soit $\approx 5{,}8$ millions de visites par jour) :

\begin{center}
\small
\setlength{\tabcolsep}{4pt}
\renewcommand{\arraystretch}{1.25}
\begin{tabular}{@{}p{0.28\linewidth}p{0.34\linewidth}p{0.32\linewidth}@{}}
  \toprule
  \textbf{Critère sur 12 mois} & \textbf{Datamart Updated-Records} \newline (\texttt{FAIT\_CYCLE\_ABONNEMENT}) & \textbf{Datamart Snapshot} \newline (\texttt{FAITS\_CONSULTATION\_QUOTIDIENNE}) \\
  \midrule
  \textbf{Hypothèse de calcul} & $670\,000$ contrats initiaux $+ 20\,\%$ de nouveaux contrats sur l'année ($+134\,000$). Les contrats résiliés restent en base ($+0$ ligne). & $\approx 15\,000$ articles consultés/jour (nouveautés + archives) $\times 150$ combinaisons actives (\emph{source $\times$ support $\times$ profil}) $= 2{,}25\,\text{M}$ lignes/jour. \\
  \addlinespace
  \textbf{Nombre de lignes à 1 an} & \textbf{$\approx 804\,000$ lignes} & \textbf{$\approx 820$ millions de lignes} \\
  \addlinespace
  \textbf{Opérations d'écriture} & $\approx 6{,}6$ millions de mises à jour d'étapes/paiements + rafraîchissement nocturne des actifs ($\approx 255\,\text{M}$ écritures/an). & Chargement par ajout (\emph{append}) de $2{,}25$ millions de nouvelles lignes chaque nuit. \\
  \addlinespace
  \textbf{Compatibilité Excel} \newline {\footnotesize (limite : $1\,048\,576$ lignes)} & Tient la 1\textsuperscript{re} année en nombre de lignes, mais dépasse la limite dès la 3\textsuperscript{e} année ($\approx 1{,}07\,\text{M}$ lignes). & \textbf{Impossible dès le 1\textsuperscript{er} jour} ($2{,}25\,\text{M}$ lignes/jour, soit $780$ fois la limite d'Excel sur un an). \\
  \bottomrule
\end{tabular}
\end{center}

\paragraph{Pourquoi un fichier Excel ne suffit dans aucun des deux cas ?}
Pour le \textbf{Snapshot des consultations}, la volumétrie annuelle ($\approx 820$ millions de lignes, qui peut être réduite d'un facteur 5 en simplifiant la finesse technique de \texttt{DIM\_SUPPORT}) impose d'emblée un SGBD orienté entrepôt de données avec partitionnement temporel et indexation adaptée.

Pour le modèle \textbf{Updated-Records}, bien que ses $804\,000$ lignes tiennent théoriquement dans une feuille Excel la première année, l'usage d'un tableur est exclu pour trois raisons :
\begin{enumerate}
  \item \textbf{Intensité des mises à jour :} Réécrire automatiquement des centaines de milliers de lignes chaque nuit par clé de contrat (\texttt{MERGE / UPDATE}) est ingérable dans un tableur.
  \item \textbf{Intégrité et partage des dimensions :} Ce modèle comporte 13 clés étrangères (dont 6 rôles de date) et partage ses dimensions conformes avec le datamart détaillé de la Partie 2.
  \item \textbf{Sécurité et RGPD :} La gestion de données contractuelles sur près d'un million d'abonnés exige un contrôle d'accès granulaire, une traçabilité et une isolation que seul un SGBD garantit.
\end{enumerate}

% ============================================================
%  PARTIE 4 -- Techniques avancées de modélisation
% ============================================================

\newpage
\section{Techniques avancées de modélisation}

Cette section approfondit l'architecture de l'entrepôt de données du \emph{Monde} en mettant en œuvre trois techniques avancées de modélisation dimensionnelle : la résolution d'une association plusieurs-à-plusieurs ($N:M$) par une table-pont, la gestion de l'évolution des attributs de la dimension la plus volumineuse (\texttt{DIM\_LECTEUR}), et la mise en place d'un schéma de partitionnement hybride pour réduire les lectures disque lors des requêtes analytiques.

% ------------------------------------------------------------
\subsection{Modélisation d'une relation $N:M$ par table-pont}
% ------------------------------------------------------------

\paragraph{Problématique métier.}
Dans le modèle \emph{Snapshot} de la Partie 3, la dimension \texttt{DIM\_ARTICLE} ne retient qu'un seul \texttt{auteur\_principal}. Or, dans un grand quotidien comme \emph{Le Monde}, les enquêtes approfondies, les couvertures électorales ou les dossiers multimédias sont fréquemment \textbf{co-rédigés par plusieurs journalistes} (par exemple un reporter sur le terrain et un analyste en rédaction). Réciproquement, un journaliste rédige de nombreux articles au fil de l'année. Il existe donc une relation plusieurs-à-plusieurs ($N:M$) entre \texttt{DIM\_ARTICLE} et les journalistes.

Relier directement une table \texttt{DIM\_AUTEUR} à la table de faits sans précaution provoquerait un problème de \textbf{double comptage} : un article totalisant $10\,000$ vues et signé par trois journalistes verrait son audience triplée lors d'une agrégation globale.

\paragraph{Conception de la table-pont.}
Pour modéliser proprement cette association sans fausser les mesures additives de \texttt{FAITS\_CONSULTATION\_QUOTIDIENNE}, nous introduisons une dimension \texttt{DIM\_AUTEUR} et une table-pont intermédiaire \texttt{PONT\_ARTICLE\_AUTEUR}. Cette table porte un attribut indispensable : le \textbf{facteur de pondération} (\texttt{facteur\_ponderation}), égal à $1/k$ pour un article rédigé par $k$ auteurs, garantissant que la somme des poids pour un article donné vaille toujours $1{,}0$.

\begin{figure}[htbp]
\centering
\begin{tikzpicture}[
  font=\footnotesize,
  dim/.style={
    rectangle split, rectangle split parts=2,
    rectangle split part fill={lightblue,white},
    draw=black!55, line width=0.6pt, rounded corners=2pt,
    text width=4.3cm, align=left, inner sep=4pt
  },
  bridge/.style={
    rectangle split, rectangle split parts=2,
    rectangle split part fill={lightgreen,white},
    draw=codegreen, line width=1.1pt, rounded corners=2pt,
    text width=4.6cm, align=left, inner sep=4pt
  },
  fact/.style={
    rectangle split, rectangle split parts=3,
    rectangle split part fill={lightorange,white,lightgreen},
    draw=codeblue, line width=1.2pt, rounded corners=2pt,
    text width=4.6cm, align=left, inner sep=4pt
  },
  lien/.style={draw=black!55, line width=0.8pt},
  card/.style={font=\scriptsize\bfseries, text=black!75}
]

\node[fact] (FAIT) at (-5.6,0) {
  \textbf{FAITS\_CONSULTATION\_QUOT.}\\
  {\scriptsize\itshape Table de faits Snapshot}
  \nodepart{two}
  date\_sk (FK)\\
  \textbf{article\_sk (FK)}\\
  source\_sk, support\_sk, profil\_sk
  \nodepart{three}
  nb\_vues\_totales \hfill [A]\\
  temps\_lecture\_total\_sec \hfill [A]
};

\node[dim] (ART) at (0,0) {
  \textbf{DIM\_ARTICLE}\\
  {\scriptsize\itshape Contenu éditorial}
  \nodepart{two}
  \textbf{article\_sk (PK)}\\
  id\_article\_source\\
  titre\_article\\
  rubrique\_principale\\
  type\_acces
};

\node[bridge] (PONT) at (5.6,0) {
  \textbf{PONT\_ARTICLE\_AUTEUR}\\
  {\scriptsize\itshape Table-pont ($N:M$)}
  \nodepart{two}
  \textbf{article\_sk (PK, FK)}\\
  \textbf{auteur\_sk (PK, FK)}\\
  \textbf{facteur\_ponderation} ($0 < w \le 1$)\\
  role\_redaction (Principal, Co-auteur)\\
  ordre\_signature (1, 2, 3...)
};

\node[dim, anchor=north] (AUT) at ($(PONT.south)+(0,-1.1)$) {
  \textbf{DIM\_AUTEUR}\\
  {\scriptsize\itshape Journalistes et contributeurs}
  \nodepart{two}
  \textbf{auteur\_sk (PK)}\\
  id\_journaliste\_source\\
  nom\_complet\\
  service\_redaction\\
  statut (Titulaire, Pigiste, Chroniqueur)\\
  date\_entree\_redaction
};

\draw[lien] (FAIT.east) -- (ART.west)
  node[card, pos=0.15, above] {$N$}
  node[card, pos=0.85, above] {$1$};

\draw[lien] (ART.east) -- (PONT.west)
  node[card, pos=0.15, above] {$1$}
  node[card, pos=0.85, above] {$N$};

\draw[lien] (PONT.south) -- (AUT.north)
  node[card, pos=0.2, right] {$N$}
  node[card, pos=0.8, right] {$1$};

\end{tikzpicture}
\caption{Intégration de la table-pont \texttt{PONT\_ARTICLE\_AUTEUR} pour modéliser la relation $N:M$ entre articles et journalistes.}
\label{fig:table-pont-auteur}
\end{figure}

\paragraph{Requête analytique exploitant la table-pont.}
La requête suivante mesure l'audience générée par chaque service rédactionnel sur les articles \emph{Premium} en 2026. La multiplication par \texttt{facteur\_ponderation} garantit que la somme des vues attribuées à l'ensemble des journalistes reste strictement égale au total réel des vues du journal :

\begin{lstlisting}[language=SQL]
SELECT au.service_redaction,
       au.nom_complet,
       COUNT(DISTINCT f.article_sk) AS nb_articles_signes,
       ROUND(SUM(f.nb_vues_totales * p.facteur_ponderation)) AS vues_ponderees,
       ROUND(SUM(f.temps_lecture_total_sec * p.facteur_ponderation) / 60.0) AS minutes_lues_ponderees
FROM faits_consultation_quotidienne f
JOIN dim_date d            ON f.date_sk = d.date_sk
JOIN dim_article a         ON f.article_sk = a.article_sk
JOIN pont_article_auteur p ON a.article_sk = p.article_sk
JOIN dim_auteur au         ON p.auteur_sk = au.auteur_sk
WHERE d.annee = 2026
  AND a.type_acces = 'Premium'
GROUP BY au.service_redaction, au.nom_complet
ORDER BY vues_ponderees DESC;
\end{lstlisting}

\noindent\textit{Cet exemple illustre le calcul de la contribution éditoriale pondérée de chaque journaliste sans double comptage en cas de co-signature.}

% ------------------------------------------------------------
\subsection{Dimension la plus volumineuse : attributs statiques et dynamiques}
% ------------------------------------------------------------

\paragraph{Identification de la dimension la plus volumineuse.}
Dans notre entrepôt de données, la dimension la plus volumineuse est \textbf{\texttt{DIM\_LECTEUR}} (partagée entre le modèle détaillé de la Partie 2 et le modèle \emph{Updated-Records} de la Partie 3). Alors que les tables de paramétrage (\texttt{DIM\_OFFRE}, \texttt{DIM\_CANAL}, \texttt{DIM\_PROMOTION}, \texttt{DIM\_STATUT\_CONTRAT}, \texttt{DIM\_SUPPORT}) ne comptent que quelques dizaines ou centaines de lignes, que \texttt{DIM\_DATE} contient $\approx 3\,650$ lignes sur 10 ans et \texttt{DIM\_GEOGRAPHIE} $\approx 35\,000$ communes, \texttt{DIM\_LECTEUR} regroupe les $670\,000$ abonnés numériques actifs, les anciens abonnés résiliés ainsi que plusieurs millions de lecteurs inscrits gratuitement. Son historisation en SCD Type 2 multiplie en outre le nombre de lignes physiques à chaque évolution de profil.

\paragraph{Classification des attributs de \texttt{DIM\_LECTEUR}.}
Nous distinguons deux catégories d'attributs selon leur évolutivité au cours du cycle de vie d'un compte :
\begin{itemize}
  \item \textbf{Attributs statiques (immuables) :} \texttt{id\_compte\_source} (identifiant unique du lecteur dans le système opérationnel CRM) et \texttt{date\_creation\_compte} (date de première inscription). Ces attributs caractérisent la création du compte et ne sont jamais modifiés.
  \item \textbf{Attributs dynamiques (évolutifs) :} \texttt{est\_etudiant}, \texttt{categorie\_socio\_pro}, \texttt{tranche\_age}, \texttt{niveau\_engagement}, \texttt{rubriques\_preferees}, \texttt{temps\_visite\_jour} et \texttt{temps\_visite\_semaine}. Ces attributs évoluent avec la situation personnelle ou les habitudes de lecture de l'utilisateur.
\end{itemize}

% ------------------------------------------------------------
\subsection{Exemples de mise à jour et stratégies SCD (Types 1, 2 et 3)}
% ------------------------------------------------------------

Tous les attributs dynamiques n'ont pas les mêmes contraintes de traçabilité historique. Le tableau ci-dessous présente un exemple concret de mise à jour pour chaque attribut et justifie la stratégie SCD (\emph{Slowly Changing Dimension}) la plus adaptée :

\begin{center}
\footnotesize
\setlength{\tabcolsep}{4pt}
\renewcommand{\arraystretch}{1.25}
\begin{tabular}{@{}p{0.21\linewidth}p{0.31\linewidth}p{0.14\linewidth}p{0.28\linewidth}@{}}
  \toprule
  \textbf{Attribut dynamique} & \textbf{Exemple de mise à jour} & \textbf{Stratégie SCD} & \textbf{Justification métier} \\
  \midrule
  \texttt{est\_etudiant} & Passe de \texttt{Oui} à \texttt{Non} à la fin des études (expiration du justificatif annuel). & \textbf{Type 2} \newline (Nouvelle ligne) & Indispensable pour conserver l'historique exact : les paiements passés restent liés au tarif étudiant, les nouveaux au plein tarif. \\
  \addlinespace
  \texttt{categorie\_socio\_pro} & Un lecteur passe de \texttt{Étudiant} à \texttt{Cadre} lors de son entrée dans la vie active. & \textbf{Type 2} \newline (Nouvelle ligne) & Permet d'analyser fidèlement le taux de rétention et le revenu moyen par CSP au moment de chaque événement. \\
  \addlinespace
  \texttt{tranche\_age} & Un lecteur atteint 25 ans et bascule de \texttt{18--24 ans} à \texttt{25--34 ans}. & \textbf{Type 3} \newline (Colonne précédente) & Ajout d'un attribut \texttt{tranche\_age\_precedente} : évite de cloner des millions de lignes chaque année tout en gardant la trace de la tranche d'origine. \\
  \addlinespace
  \texttt{rubriques\_preferees} & Le centre d'intérêt dominant glisse de \texttt{Sport} vers \texttt{Économie}. & \textbf{Type 3} \newline (Colonne précédente) & L'attribut \texttt{rubrique\_pref\_initiale} conserve la rubrique qui a déclenché l'abonnement, tandis que \texttt{rubriques\_preferees} reflète l'usage actuel. \\
  \addlinespace
  \texttt{niveau\_engagement} & Passe de \texttt{Quotidien} à \texttt{Occasionnel} suite à une baisse de fréquence de lecture. & \textbf{Type 3} \newline (Colonne précédente) & Ajout d'un attribut \texttt{niveau\_engagement\_prec} : détecte immédiatement un décrochage récent pour déclencher une action de rétention sans dupliquer la ligne. \\
  \addlinespace
  \texttt{temps\_visite\_jour} \newline \texttt{temps\_visite\_semaine} & La tranche de visite passe de \texttt{5--15 min} à \texttt{15--30 min} lors du recalcul mensuel. & \textbf{Type 1} \newline (Écrasement) & Métriques comportementales lissées recalculées chaque mois ; les historiser en Type 2 provoquerait une explosion combinatoire de lignes. \\
  \bottomrule
\end{tabular}
\end{center}

% ------------------------------------------------------------
\subsection{Schéma de partitionnement hybride sur \texttt{DIM\_LECTEUR}}
% ------------------------------------------------------------

\paragraph{Architecture du partitionnement hybride.}
Afin d'éviter que les mises à jour comportementales mensuelles ne verrouillent et ne réécrivent inutilement les attributs socio-démographiques sur plusieurs millions de lignes, nous mettons en place un \textbf{partitionnement hybride} combinant un découpage vertical et un découpage horizontal :

\begin{enumerate}
  \item \textbf{Partitionnement vertical (par colonnes, selon l'évolutivité) :} La table monolithique \texttt{DIM\_LECTEUR} est scindée en deux tables physiques plus étroites reliées par la clé primaire \texttt{lecteur\_sk} (et réunifiées logiquement par une vue SQL \texttt{dim\_lecteur}) :
  \begin{itemize}
    \item \texttt{DIM\_LECTEUR\_PROFIL} regroupe les attributs statiques et socio-démographiques (SCD 2 et 3) : \texttt{id\_compte\_source}, \texttt{date\_creation\_compte}, \texttt{est\_etudiant}, \texttt{categorie\_socio\_pro}, \texttt{tranche\_age}, \texttt{tranche\_age\_prec}, \texttt{rubriques\_preferees}, \texttt{rubrique\_pref\_init}, \texttt{debut\_validite} et \texttt{fin\_validite}.
    \item \texttt{DIM\_LECTEUR\_ACTIVITE} isole les attributs comportementaux mis à jour mensuellement (SCD 1 et 3) : \texttt{niveau\_engagement}, \texttt{niveau\_engagement\_prec}, \texttt{temps\_visite\_jour} et \texttt{temps\_visite\_semaine}.
  \end{itemize}

  \item \textbf{Partitionnement horizontal (par lignes, selon un critère analytique) :} Chacune des deux tables verticales est découpée physiquement en sous-partitions sur le disque selon un attribut de filtrage très fréquent dans les requêtes décisionnelles :
  \begin{itemize}
    \item \texttt{DIM\_LECTEUR\_PROFIL} est partitionnée par liste (\texttt{PARTITION BY LIST}) sur \textbf{\texttt{est\_etudiant}} (\texttt{'Oui'} \emph{vs} \texttt{'Non'}), isolant la population étudiante à forte rotation annuelle et tarification spécifique.
    \item \texttt{DIM\_LECTEUR\_ACTIVITE} est partitionnée par liste sur \textbf{\texttt{niveau\_engagement}} (\texttt{'Quotidien'}, \texttt{'Hebdomadaire'}, \texttt{'Occasionnel'}, \texttt{'Inactif'}), critère central pour les campagnes de rétention et de montée en gamme.
  \end{itemize}
\end{enumerate}

\begin{lstlisting}[language=SQL]
-- 1. Partitionnement vertical (attributs statiques et SCD 2/3)
--    + Partitionnement horizontal par liste sur le statut etudiant
CREATE TABLE dim_lecteur_profil (
    lecteur_sk            BIGINT NOT NULL,
    id_compte_source      VARCHAR(32) NOT NULL,
    date_creation_compte  DATE NOT NULL,
    est_etudiant          VARCHAR(3) NOT NULL, -- 'Oui' ou 'Non'
    categorie_socio_pro   VARCHAR(50),
    tranche_age           VARCHAR(20),
    tranche_age_prec      VARCHAR(20),         -- SCD Type 3
    rubriques_preferees   VARCHAR(100),
    rubrique_pref_init    VARCHAR(100),        -- SCD Type 3
    debut_validite        DATE NOT NULL,
    fin_validite          DATE NOT NULL,
    PRIMARY KEY (lecteur_sk, est_etudiant)
) PARTITION BY LIST (est_etudiant);

CREATE TABLE dim_lecteur_profil_etudiant PARTITION OF dim_lecteur_profil FOR VALUES IN ('Oui');
CREATE TABLE dim_lecteur_profil_standard PARTITION OF dim_lecteur_profil FOR VALUES IN ('Non');

-- 2. Partitionnement vertical (attributs dynamiques comportementaux SCD 1/3)
--    + Partitionnement horizontal par niveau d'engagement
CREATE TABLE dim_lecteur_activite (
    lecteur_sk             BIGINT NOT NULL,
    niveau_engagement      VARCHAR(20) NOT NULL,
    niveau_engagement_prec VARCHAR(20),        -- SCD Type 3
    temps_visite_jour      VARCHAR(20),
    temps_visite_semaine   VARCHAR(20),
    PRIMARY KEY (lecteur_sk, niveau_engagement)
) PARTITION BY LIST (niveau_engagement);

CREATE TABLE dim_lecteur_act_quotidien   PARTITION OF dim_lecteur_activite FOR VALUES IN ('Quotidien');
CREATE TABLE dim_lecteur_act_hebdo       PARTITION OF dim_lecteur_activite FOR VALUES IN ('Hebdomadaire');
CREATE TABLE dim_lecteur_act_occasionnel PARTITION OF dim_lecteur_activite FOR VALUES IN ('Occasionnel');
CREATE TABLE dim_lecteur_act_inactif     PARTITION OF dim_lecteur_activite FOR VALUES IN ('Inactif');
\end{lstlisting}

\paragraph{Pourquoi ce partitionnement réduit-il les lectures de données (I/O) ?}
Ce double découpage physique diminue fortement le volume de blocs lus sur le disque lors de l'évaluation des requêtes :
\begin{itemize}
  \item \textbf{Densité accrue des blocs disque (effet du partitionnement vertical) :} Dans un stockage orienté lignes, chaque bloc disque (ex. $8\,\text{Ko}$) contient un nombre fixe d'octets. En extrayant les colonnes d'activité dans \texttt{dim\_lecteur\_activite}, chaque ligne ne pèse que quelques octets (5 colonnes étroites au lieu de 13). Le SGBD stocke ainsi \textbf{3 à 4 fois plus d'enregistrements par bloc disque}, réduisant d'autant le nombre de lectures physiques nécessaires pour scanner l'activité des lecteurs. De même, les mises à jour mensuelles SCD 1/3 ne réécrivent que ces blocs légers sans toucher aux attributs socio-démographiques.
  \item \textbf{Élagage de partitions (\emph{Partition Pruning}, effet du partitionnement horizontal) :} Lorsqu'une requête filtre sur le critère de partitionnement (par exemple \texttt{WHERE niveau\_engagement = 'Inactif'} pour cibler une campagne anti-résiliation, ou \texttt{WHERE est\_etudiant = 'Oui'}), l'optimiseur SQL identifie immédiatement la sous-table physique concernée et \textbf{ignore sans les lire} toutes les autres partitions sur le disque. Au lieu de parcourir 100\,\% des millions de lignes de la dimension, le moteur ne charge que le sous-ensemble utile (environ 10 à 20\,\% du volume total), divisant proportionnellement les entrées/sorties disque (I/O) et le temps d'exécution.
\end{itemize}

\section{L’implémentation et le requêtage, les Vues Matérialisées, et les Indexes Bitmap}

Le modèle détaillé de la Partie~2 a été implémenté sous Oracle~19c, sur le
serveur de l'université. Le script \texttt{modele\_detaille.sql} crée les huit
tables de dimensions et la table des faits \texttt{FAIT\_EVENEMENT\_ABONNEMENT},
avec leurs clés primaires, clés étrangères et contraintes de contrôle, puis
insère les données de l'instance présentée en Partie~2.
Les dimensions partagées sont exposées par des vues virtuelles :
\texttt{v\_date\_evenement} et \texttt{v\_date\_souscription} donnent chacune
un rôle à \texttt{DIM\_DATE}, et \texttt{v\_lecteur\_courant} ne garde que la
version actuelle de chaque lecteur.
Le fichier \texttt{requetes\_analytiques.sql} regroupe cinq requêtes posées sur
ce même modèle en étoile. Elles répondent aux traitements de l'Action~1 et
utilisent notamment les opérateurs \texttt{ROLLUP} et \texttt{CUBE}.
Le fichier \texttt{vues\_materialisees.sql} définit deux vues matérialisées,
choisies à l'aide du treillis d'agrégation : la première sert les requêtes
sur les flux d'abonnés (R1 et R2), la seconde celles sur l'activité
commerciale (R3, R4 et R5).
Enfin, le fichier \texttt{bitmap\_join.sql} crée un index bitmap de jointure
sur le type d'événement et l'année, qui optimise les requêtes R3 et R4.
Les plans d'exécution obtenus avec \texttt{EXPLAIN PLAN} confirment son
utilisation.
Tous ces scripts sont fournis avec le rendu. Chaque étape y est expliquée en
commentaire, ainsi que les choix de modélisation. Pour vérifier les résultats, le plus intéressant et de lancer les fichiers dans l'ordre demandé dans le sujet pour voir sur le terminal ce que ça renvoie.

\newpage

\section*{Conclusion}
\addcontentsline{toc}{section}{Conclusion}

Nous en avons terminé avec l'implémentation : elle fonctionne et permet de
retracer toute notre logique métier depuis le début. Nous avons d'abord vu
que \emph{Le Monde} est un journal qui évolue au fil des années. Son modèle
économique repose désormais sur les abonnements numériques plutôt que sur
la vente papier ou la publicité, car l'information se consomme aujourd'hui
différemment. Cette transition oblige le journal à comprendre précisément
qui s'abonne, pourquoi, et qui risque de partir. C'est le rôle de
l'entrepôt de données que nous avons conçu.

Le modèle détaillé, centré sur les événements d'abonnement, donne une vue
d'ensemble de l'activité du journal. Ses dimensions permettent d'analyser
chaque entrée et chaque sortie d'abonné selon le profil du lecteur, l'offre,
le canal, la promotion, le lieu et le moment. Les deux modèles moins
détaillés le complètent sous d'autres angles. Le \emph{snapshot} résume
chaque jour la consultation des articles, ce qui rend l'audience analysable
malgré des millions de clics quotidiens. L'\emph{updated-records} suit le
cycle de vie de chaque contrat, de la souscription à la résiliation, pour
mesurer la fidélité et la rentabilité des promotions. Les techniques
avancées (table-pont, historisation des lecteurs, partitionnement) et
l'implémentation sous Oracle, avec ses vues, vues matérialisées et index
bitmap de jointure, montrent que ces modèles peuvent être interrogés
efficacement sur de grands volumes.

Notre travail a toutefois des limites. Les données de test sont fictives,
et plusieurs volumétries reposent sur des hypothèses. Surtout, aucun de nos
modèles ne relie encore un abonnement à l'article qui l'a déclenché, alors
que c'était l'un des premiers besoins identifiés. Ajouter cette information,
par exemple avec une clé vers l'article déclencheur dans la table des faits,
serait la prochaine étape pour que l'entrepôt réponde pleinement à la
question centrale du journal : quels contenus donnent envie de s'abonner ?

% ============================================================
%  ANNEXES -- Tables des exemples d'instance
% ============================================================

\newpage
\appendix

\section{Instance du modèle détaillé}
\label{annexe:instance-detaille}

\begin{center}
\scriptsize
\setlength{\tabcolsep}{3pt}
\renewcommand{\arraystretch}{1.1}

% ---------- DIM_DATE | DIM_HEURE ----------
\begin{minipage}[t]{0.49\linewidth}\centering
\textbf{DIM\_DATE}\\[2pt]
\begin{tabular}{@{}llllll@{}}
  \toprule
  date\_sk & date & jour & WE & mois & année \\
  \midrule
  20240110 & 2024-01-10 & Mer. & Non & Janv. & 2024 \\
  20260110 & 2026-01-10 & Sam. & Oui & Janv. & 2026 \\
  20260115 & 2026-01-15 & Jeu. & Non & Janv. & 2026 \\
  20260215 & 2026-02-15 & Dim. & Oui & Févr. & 2026 \\
  \bottomrule
\end{tabular}
\end{minipage}%
\hfill%
\begin{minipage}[t]{0.49\linewidth}\centering
\textbf{DIM\_HEURE}\\[2pt]
\begin{tabular}{@{}lllll@{}}
  \toprule
  heure\_sk & heure & tranche & pointe & après NL \\
  \midrule
  0005 & 00:05 & Nuit & Non & Non \\
  1240 & 12:40 & Midi & Oui & Non \\
  2130 & 21:30 & Soir & Non & Non \\
  \bottomrule
\end{tabular}
\end{minipage}

\medskip

% ---------- DIM_LECTEUR | DIM_OFFRE ----------
\begin{minipage}[t]{0.49\linewidth}\centering
\textbf{DIM\_LECTEUR}\\[2pt]
\begin{tabular}{@{}llllll@{}}
  \toprule
  sk & âge & CSP & étud. & visite/jour & rubrique préf. \\
  \midrule
  1001 & 18--24 & Étudiant & Oui & 5--15 min  & International \\
  1002 & 55--64 & Cadre    & Non & 15--30 min & Économie \\
  \bottomrule
\end{tabular}
\end{minipage}%
\hfill%
\begin{minipage}[t]{0.49\linewidth}\centering
\textbf{DIM\_OFFRE}\\[2pt]
\begin{tabular}{@{}llllll@{}}
  \toprule
  sk & code & premium & période & engag. & prix TTC \\
  \midrule
  12 & NUM-ESS-M & Non & Mensuelle & 0  & 11,99\,€ \\
  20 & NUM-INT-A & Oui & Annuelle  & 12 & 199,00\,€ \\
  \bottomrule
\end{tabular}
\end{minipage}

\medskip

% ---------- DIM_CANAL | DIM_PROMOTION ----------
\begin{minipage}[t]{0.49\linewidth}\centering
\textbf{DIM\_CANAL}\\[2pt]
\begin{tabular}{@{}llll@{}}
  \toprule
  sk & libellé & famille & appareil \\
  \midrule
  5 & Paywall article     & Paywall    & Mobile \\
  8 & Newsletter matinale & Newsletter & Ordinateur \\
  \bottomrule
\end{tabular}
\end{minipage}%
\hfill%
\begin{minipage}[t]{0.49\linewidth}\centering
\textbf{DIM\_PROMOTION}\\[2pt]
\begin{tabular}{@{}lllll@{}}
  \toprule
  sk & code & type & durée & prix promo \\
  \midrule
  0 & --- & Aucune promotion & --- & --- \\
  3 & PREMIER1E & Prix fixe & 1 mois & 1,00\,€ \\
  \bottomrule
\end{tabular}
\end{minipage}

\medskip

% ---------- DIM_GEOGRAPHIE | DIM_TYPE_EVENEMENT ----------
\begin{minipage}[t]{0.49\linewidth}\centering
\textbf{DIM\_GEOGRAPHIE}\\[2pt]
\begin{tabular}{@{}llll@{}}
  \toprule
  sk & commune & département & région \\
  \midrule
  690 & Lyon   & Rhône (69)           & Auvergne-Rhône-Alpes \\
  350 & Rennes & Ille-et-Vilaine (35) & Bretagne \\
  \bottomrule
\end{tabular}
\end{minipage}%
\hfill%
\begin{minipage}[t]{0.49\linewidth}\centering
\textbf{DIM\_TYPE\_EVENEMENT}\\[2pt]
\begin{tabular}{@{}llllll@{}}
  \toprule
  sk & type & sens & motif & initiateur & paiement \\
  \midrule
  1 & Souscription   & Entrée   & --- & Abonné  & Accepté \\
  2 & Renouvellement & Maintien & --- & Système & Accepté \\
  3 & Résiliation    & Sortie   & Prix & Abonné & --- \\
  \bottomrule
\end{tabular}
\end{minipage}

\bigskip

% ---------- FAIT : clés ----------
\textbf{FAIT\_EVENEMENT\_ABONNEMENT -- clés et dimensions dégénérées}\\[2pt]
\begin{tabular}{@{}llllllllllll@{}}
  \toprule
  Ligne & d\_evt & d\_sous & heure & lect. & offre & canal & promo & geo & type & contrat & transaction \\
  \midrule
  F1 & 20260115 & 20260115 & 1240 & 1001 & 12 & 5 & 3 & 690 & 1 & CT-5001 & TR-90001 \\
  F2 & 20260215 & 20260115 & 0005 & 1001 & 12 & 5 & 0 & 690 & 2 & CT-5001 & TR-90417 \\
  F3 & 20260110 & 20240110 & 2130 & 1002 & 20 & 8 & 0 & 350 & 3 & CT-1187 & --- \\
  \bottomrule
\end{tabular}

\medskip

% ---------- FAIT : mesures ----------
\textbf{FAIT\_EVENEMENT\_ABONNEMENT -- mesures}\\[2pt]
\begin{tabular}{@{}lrrrrrrrr@{}}
  \toprule
  Ligne & ttc & ht & remise & mois & delta & revenu cumulé & ancienneté & taux remise \\
  \midrule
  F1 & 1,00  & 0,98  & 10,99 & 1 & $+1$ & 1,00   & 0  & 91,7\,\% \\
  F2 & 11,99 & 11,74 & 0,00  & 1 & $0$  & 12,99  & 1  & 0,0\,\%  \\
  F3 & 0,00  & 0,00  & 0,00  & 0 & $-1$ & 398,00 & 24 & ---      \\
  \bottomrule
\end{tabular}
\end{center}

\section{Instance du modèle snapshot}
\label{annexe:instance-snapshot}

\begin{center}
\footnotesize
\setlength{\tabcolsep}{4pt}
\renewcommand{\arraystretch}{1.15}

\begin{tabular}{@{}llllll@{}}
  \toprule
  \multicolumn{6}{@{}l}{\textbf{DIM\_ARTICLE} (extrait)} \\
  \textbf{article\_sk} & \textbf{titre\_article} & \textbf{rubrique} & \textbf{type\_acces} & \textbf{format} & \textbf{longueur\_mots} \\
  \midrule
  501 & Les enjeux des élections... & Politique & Premium & Texte & 1\,500 \\
  502 & Résultats du marathon & Sport & Gratuit & Vidéo & 300 \\
  \bottomrule
\end{tabular}

\medskip

\begin{tabular}{@{}lllll@{}}
  \toprule
  \multicolumn{5}{@{}l}{\textbf{DIM\_SOURCE\_TRAFIC} (extrait)} \\
  \textbf{source\_sk} & \textbf{nom\_source} & \textbf{categorie\_source} & \textbf{medium\_acquisition} & \textbf{est\_trafic\_interne} \\
  \midrule
  11 & Google News & Moteur de recherche & organic & Non \\
  12 & Page d'accueil & Site web & direct & Oui \\
  \bottomrule
\end{tabular}

\medskip

\begin{tabular}{@{}llll@{\hspace{0.5cm}}llll@{}}
  \toprule
  \multicolumn{4}{@{}l}{\textbf{DIM\_SUPPORT} (extrait)} & \multicolumn{4}{@{}l}{\textbf{DIM\_PROFIL\_LECTEUR} (extrait)} \\
  \textbf{sup\_sk} & \textbf{appareil} & \textbf{OS} & \textbf{application} & \textbf{profil\_sk} & \textbf{statut\_abo} & \textbf{engagement} & \textbf{tranche\_age} \\
  \midrule
  81 & Mobile & iOS & App Native & 201 & Abonné Intégral & Quotidien & 35--49 ans \\
  82 & Ordinateur & Windows & Nav. Web & 202 & Visiteur Gratuit & Occasionnel & 18--24 ans \\
  \bottomrule
\end{tabular}

\medskip

\begin{tabular}{@{}rrrrrrrrrrr@{}}
  \toprule
  \multicolumn{11}{@{}l}{\textbf{FAITS\_CONSULTATION\_QUOTIDIENNE}} \\
  \textbf{date\_sk} & \textbf{art\_sk} & \textbf{src\_sk} & \textbf{sup\_sk} & \textbf{seg\_sk} & \textbf{vues} & \textbf{compl.} & \textbf{tps\_tot (s)} & \textbf{comm.} & \textbf{tps\_moy} & \textbf{tx\_compl} \\
  \midrule
  20261029 & 501 & 12 & 81 & 201 & 5\,000  & 4\,000 & 900\,000 & 45 & 180\,s & 80\,\% \\
  20261029 & 502 & 11 & 82 & 202 & 12\,000 & 3\,000 & 720\,000 & 10 & 60\,s  & 25\,\% \\
  \bottomrule
\end{tabular}
\end{center}

\section{Instance du modèle updated-records}
\label{annexe:instance-updated}

\begin{center}
\small
\setlength{\tabcolsep}{3.5pt}
\renewcommand{\arraystretch}{1.15}
\begin{tabular}{@{}llllll@{}}
  \toprule
  \multicolumn{6}{@{}l}{\textbf{DIM\_STATUT\_CONTRAT} (extrait)} \\
  \textbf{statut\_sk} & \textbf{statut} & \textbf{etape} & \textbf{est\_actif} & \textbf{plein\_tarif} & \textbf{type / motif / initiateur\_fin} \\
  \midrule
  1 & En promotion & 1 & Oui & Non & --- \\
  2 & Plein tarif  & 2 & Oui & Oui & --- \\
  5 & Résilié      & 5 & Non & Non & Volontaire / Prix / Abonné \\
  \bottomrule
\end{tabular}

\bigskip

\setlength{\tabcolsep}{2.5pt}
\begin{tabular}{@{}llllllllllll@{}}
  \toprule
  \multicolumn{12}{@{}l}{\textbf{FAIT\_CYCLE\_ABONNEMENT -- Clés étrangères et jalons}} \\
  \textbf{État} & \textbf{Contrat} & \textbf{d\_sous} & \textbf{d\_fin\_pr} & \textbf{d\_plein} & \textbf{d\_dem\_r} & \textbf{d\_fin} & \textbf{d\_maj} & \textbf{lect} & \textbf{off\_i/a} & \textbf{pro} & \textbf{stat} \\
  \midrule
  $E_1$ & CT-5001 & 20260115 & 99991231 & 99991231 & 99991231 & 99991231 & 20260115 & 1001 & 12 / 12 & 3 & 1 \\
  $E_2$ & CT-5001 & 20260115 & \textbf{20260215} & \textbf{20260215} & 99991231 & 99991231 & \textbf{20260215} & 1001 & 12 / 12 & 3 & \textbf{2} \\
  $E_3$ & CT-5001 & 20260115 & 20260215 & 20260215 & 99991231 & 99991231 & \textbf{20260315} & 1001 & 12 / 12 & 3 & 2 \\
  \midrule
  Final & CT-1187 & 20240110 & 99991231 & 20240110 & 20251201 & 20260110 & 20260110 & 1002 & 20 / 20 & 0 & 5 \\
  \bottomrule
\end{tabular}

\bigskip

\setlength{\tabcolsep}{2.2pt}
\begin{tabular}{@{}lccccccccccccc@{}}
  \toprule
  \multicolumn{14}{@{}l}{\textbf{FAIT\_CYCLE\_ABONNEMENT -- Mesures correspondantes}} \\
  \textbf{État} & \textbf{M1} & \textbf{M2} & \textbf{M3} & \textbf{M4} & \textbf{M5} & \textbf{M6} & \textbf{M7} & \textbf{M8} & \textbf{M9} & \textbf{M10} & \textbf{M11} & \textbf{M12} & \textbf{M13} \\
  \midrule
  $E_1$ & 1 & 1,00  & 10,99 & 1  & 0 & 0 & 0 & 0 & --- & 0   & --- & 1,00  & 0 \\
  $E_2$ & 1 & \textbf{12,99} & 10,99 & \textbf{2}  & 0 & 0 & \textbf{1} & 0 & \textbf{31}  & \textbf{31}  & --- & \textbf{11,99} & 0 \\
  $E_3$ & 1 & \textbf{24,98} & 10,99 & \textbf{3}  & 0 & 0 & 1 & 0 & 31  & \textbf{59}  & --- & 11,99 & 0 \\
  \midrule
  Final & 1 & 398,00 & 0,00 & 24 & 0 & 0 & 1 & 1 & 0   & 731 & 40  & 0,00  & 0 \\
  \bottomrule
\end{tabular}
\end{center}

\listoffigures

\end{document}
